
%% ============================================================
%%  LỜI CẢM ƠN
%% ============================================================
\chapter*{Lời cảm ơn}
\addcontentsline{toc}{chapter}{Lời cảm ơn}
\markboth{LỜI CẢM ƠN}{}

Em xin bày tỏ lòng biết ơn sâu sắc đến quý thầy cô Khoa Mạng máy tính và Truyền thông, Trường Đại học Công nghệ Thông tin, Đại học Quốc gia Thành phố Hồ Chí Minh, những người đã truyền đạt nền tảng kiến thức chuyên môn, phương pháp tư duy khoa học và tinh thần làm việc nghiêm túc trong suốt quá trình học tập. Những bài giảng, những buổi thảo luận và những góp ý học thuật đã tạo nên cơ sở vững chắc để đề tài này được hình thành và phát triển theo hướng bài bản, có hệ thống và phù hợp với định hướng nghiên cứu của ngành.

Đặc biệt, em xin trân trọng cảm ơn ThS. Nguyễn Hữu Quyền, giảng viên hướng dẫn của báo cáo thực tập. Trong suốt quá trình thực hiện đề tài, từ giai đoạn xác định vấn đề nghiên cứu, xây dựng cơ sở lý thuyết, thiết kế hệ thống cho đến triển khai thực nghiệm và hoàn thiện bản thảo, sự hướng dẫn tận tình và những góp ý chi tiết của thầy đã giúp em nhìn nhận vấn đề một cách chặt chẽ và toàn diện hơn. Những buổi trao đổi chuyên môn không chỉ dừng lại ở việc chỉnh sửa nội dung, mà còn định hình phương pháp tiếp cận khoa học, cách tổ chức lập luận và cách trình bày kết quả một cách mạch lạc. Sự nghiêm túc trong học thuật và tinh thần trách nhiệm của thầy là nguồn động lực quan trọng để báo cáo này được hoàn thiện với chất lượng tốt nhất trong khả năng của em.

Bên cạnh đó, em xin gửi lời cảm ơn đến tập thể cán bộ, giảng viên và nhân viên của Khoa đã tạo điều kiện thuận lợi về cơ sở vật chất, tài liệu tham khảo và môi trường học tập.

Em cũng xin trân trọng cảm ơn Trung tâm An ninh mạng đã tạo điều kiện thực tập, cung cấp môi trường làm việc chuyên nghiệp, cơ sở hạ tầng kỹ thuật và những bài toán thực tế có giá trị, giúp em áp dụng kiến thức học thuật vào thực tiễn một cách hiệu quả.

% Em cũng xin bày tỏ lòng biết ơn đối với gia đình, những người luôn âm thầm ủng hộ và tạo điều kiện tốt nhất để việc học tập và nghiên cứu được diễn ra liên tục và ổn định. Sự động viên, tin tưởng và chia sẻ từ gia đình là chỗ dựa tinh thần quan trọng trong những giai đoạn gặp khó khăn về học thuật cũng như áp lực tiến độ.

Cuối cùng, mặc dù vẫn còn một số hạn chế và nhiều hướng có thể tiếp tục cải thiện, báo cáo đã đạt được các mục tiêu đề ra và cung cấp những kết quả ban đầu khả quan cho bài toán nghiên cứu. Những góp ý và nhận xét từ quý thầy cô sẽ là cơ sở quan trọng để tác giả tiếp tục hoàn thiện và phát triển hướng nghiên cứu này trong tương lai.

\vspace{2cm}
\begin{flushright}
\textit{TP. Hồ Chí Minh, tháng 06 năm 2026}\\[0.5cm]
% \textbf{Tác giả}
\end{flushright}

\pagebreak

%% ============================================================
%%  MỤC LỤC
%% ============================================================
\tableofcontents
\pagebreak

%% ============================================================
%%  DANH SÁCH HÌNH VẼ
%% ============================================================
\listoffigures
\pagebreak

%% ============================================================
%%  DANH SÁCH BẢNG BIỂU
%% ============================================================
\listoftables
\pagebreak

%% ============================================================
%%  DANH MỤC TỪ VIẾT TẮT
%% ============================================================
\chapter*{Danh mục từ viết tắt}
\addcontentsline{toc}{chapter}{Danh mục từ viết tắt}
\markboth{DANH MỤC TỪ VIẾT TẮT}{}

\small
\begin{longtable}{p{2.8cm} p{3.5cm} p{6.5cm}}
\hline
\textbf{Viết tắt} & \textbf{Tiếng Anh} & \textbf{Tiếng Việt} \\
\hline
\endfirsthead
\hline
\textbf{Viết tắt} & \textbf{Tiếng Anh} & \textbf{Tiếng Việt} \\
\hline
\endhead
\hline
\endfoot
\hline
\endlastfoot
APT & Advanced Persistent Threat & Mối đe dọa bền vững nâng cao \\
CNN & Convolutional Neural Network & Mạng nơ-ron tích chập \\
DL & Deep Learning & Học sâu \\
DFL & Decentralized Federated Learning & Học liên kết phân tán \\
DDoS & Distributed Denial of Service & Từ chối dịch vụ phân tán \\
EKD & Evidential Knowledge Distillation & Chắt lọc kiến thức dựa trên bằng chứng \\
EVA & Energy-based Vote Abstention & Cơ chế bỏ phiếu trắng dựa trên năng lượng \\
FD & Federated Distillation & Chắt lọc kiến thức liên kết \\
FL & Federated Learning & Học liên kết \\
GAN & Generative Adversarial Network & Mạng đối sinh \\
GNN & Graph Neural Network & Mạng nơ-ron đồ thị \\
GRU & Gated Recurrent Unit & Bộ nhớ hồi tiếp có cổng \\
IDS & Intrusion Detection System & Hệ thống phát hiện xâm nhập \\
IoT & Internet of Things & Internet vạn vật \\
LMA & Loss Maximization Attack & Tấn công tối đa hóa mất mát \\
ML & Machine Learning & Học máy \\
NIDS & Network-based IDS & IDS dựa trên mạng \\
OOD & Out-of-Distribution & Ngoài phân phối \\
RNG & Random Number Generator & Bộ sinh số ngẫu nhiên \\
RNN & Recurrent Neural Network & Mạng nơ-ron hồi tiếp \\
SDN & Software-Defined Networking & Mạng định nghĩa bằng phần mềm \\
SIEM & Security Information and Event Management & Quản lý thông tin và sự kiện bảo mật \\
SOC & Security Operations Center & Trung tâm điều hành bảo mật \\
XAI & Explainable Artificial Intelligence & Học máy khả giải \\
VM & Virtual Machine & Máy ảo \\
\end{longtable}

\pagebreak

%% ============================================================
%%  CHƯƠNG 1: TỔNG QUAN
%% ============================================================
\chapter{Tổng quan}

\section{Giới thiệu về Trung tâm An ninh mạng}
\label{sec:company}

Trung tâm An ninh mạng (Cyber Security Center) là đơn vị hoạt động trong lĩnh vực bảo mật thông tin và an toàn mạng, cung cấp các giải pháp phát hiện, ngăn chặn và ứng phó với các mối đe dọa mạng. Với đội ngũ chuyên gia giàu kinh nghiệm và hệ thống cơ sở hạ tầng kỹ thuật hiện đại, Trung tâm hướng đến việc nghiên cứu và phát triển các công nghệ tiên tiến trong lĩnh vực an ninh mạng, bao gồm hệ thống phát hiện xâm nhập (IDS), phân tích mã độc, săn tìm mối đe dọa (threat hunting) và ứng dụng trí tuệ nhân tạo trong bảo mật.

Các lĩnh vực hoạt động chính của Trung tâm bao gồm:
\begin{itemize}
    \item Nghiên cứu và phát triển các giải pháp an ninh mạng tiên tiến dựa trên trí tuệ nhân tạo và học máy.
    \item Xây dựng hệ thống giám sát an ninh mạng (SOC) và quy trình ứng phó sự cố.
    \item Phân tích mã độc và điều tra số (digital forensics).
    \item Đánh giá bảo mật hệ thống (penetration testing, vulnerability assessment).
    \item Đào tạo và nâng cao nhận thức về an ninh mạng cho tổ chức và cá nhân.
    \item Hợp tác nghiên cứu với các trường đại học, viện nghiên cứu trong và ngoài nước.
\end{itemize}

Môi trường thực tập tại Trung tâm tạo điều kiện cho sinh viên tiếp cận với các bài toán thực tế, hệ thống hạ tầng mạng quy mô lớn và những thách thức bảo mật đang nổi lên trong bối cảnh chuyển đổi số. Trong thời gian thực tập, nhóm tác giả đã được tham gia vào quá trình nghiên cứu và phát triển hệ thống phát hiện xâm nhập dựa trên học liên kết, một hướng tiếp cận mới nhằm giải quyết các vấn đề về quyền riêng tư dữ liệu, tính không đồng nhất của môi trường mạng và khả năng chống chịu trước các cuộc tấn công độc hại.

Quy trình thực tập tại Trung tâm được tổ chức theo các giai đoạn: (i) định hướng và làm quen với môi trường làm việc, (ii) nghiên cứu lý thuyết và tìm hiểu công nghệ nền tảng, (iii) thiết kế và phát triển giải pháp, (iv) thực nghiệm và đánh giá kết quả, (v) hoàn thiện báo cáo và trình bày kết quả. Trong suốt quá trình thực tập, nhóm tác giả được sự hướng dẫn trực tiếp từ các chuyên gia và cán bộ kỹ thuật của Trung tâm, đảm bảo định hướng nghiên cứu phù hợp với nhu cầu thực tế và xu hướng công nghệ mới nhất.

\section{Lý do chọn đề tài}
\label{sec:reason}

Hệ thống phát hiện xâm nhập (Intrusion Detection Systems -- IDS) là một thành phần quan trọng trong hạ tầng an ninh mạng hiện đại, cho phép xác định các hoạt động độc hại trong các mạng doanh nghiệp, đám mây, công nghiệp và Internet vạn vật (IoT). Các báo cáo gần đây cho thấy các cuộc tấn công nhắm vào IoT tiếp tục gia tăng \cite{sonicwall2024report}, trong khi các hệ thống bảo vệ nhà thông minh chặn hàng triệu sự kiện độc hại mỗi ngày \cite{netgear2024report}. Các mối đe dọa này bao gồm từ lây nhiễm mã độc, tuyển dụng botnet đến các chiến dịch tấn công từ chối dịch vụ phân tán (DDoS) quy mô lớn \cite{hackernews2026report}, làm nổi bật nhu cầu về các hệ thống phát hiện xâm nhập hiệu quả và có khả năng mở rộng.

Do đó, những tiến bộ gần đây trong Học máy (Machine Learning -- ML) và Học sâu (Deep Learning -- DL) đã cải thiện đáng kể khả năng phát hiện của IDS thông qua việc tự động học các mẫu lưu lượng phức tạp từ dữ liệu mạng quy mô lớn. Tuy nhiên, việc huấn luyện hiệu quả các IDS dựa trên ML thường yêu cầu thu thập dấu vết lưu lượng từ nhiều tổ chức, gây ra các vấn đề về quyền riêng tư, bảo mật và tuân thủ quy định. Học liên kết (Federated Learning -- FL) \cite{fedavg} đã nổi lên như một mô hình đầy hứa hẹn cho phép huấn luyện IDS cộng tác mà không chia sẻ trực tiếp dữ liệu lưu lượng riêng tư.

Mặc dù đạt được những thành công nhất định, việc triển khai thực tế FL-based IDS vẫn còn nhiều thách thức. Các mạng nơ-ron sâu hiện đại thường chứa hàng triệu tham số, dẫn đến chi phí truyền thông đáng kể trong quá trình trao đổi tham số lặp đi lặp lại. Hơn nữa, các nghiên cứu gần đây đã chỉ ra rằng các bản cập nhật mô hình có thể làm rò rỉ thông tin huấn luyện nhạy cảm thông qua các cuộc tấn công gradient inversion và reconstruction. Mặc dù các phương pháp mật mã như Secure Aggregation \cite{bonawitz2016secagg} và Homomorphic Encryption \cite{li2020deepfed} cải thiện quyền riêng tư, chúng lại tạo ra chi phí tính toán và truyền thông đáng kể.

Để giảm việc truyền tham số và hạn chế rủi ro lộ thông tin từ các bản cập nhật mô hình, các nghiên cứu gần đây đã khám phá Chắt lọc kiến thức liên kết (Federated Distillation -- FD) \cite{li2024fdsurvey}. Thay vì trao đổi tham số mô hình, FD truyền thông các biểu diễn tri thức trung gian như logits hoặc prototypes. Điều này giảm đáng kể chi phí truyền thông đồng thời hỗ trợ tự nhiên các kiến trúc mô hình không đồng nhất. Nhờ đó, FD đã trở thành một giải pháp thay thế hấp dẫn cho các môi trường hạn chế tài nguyên và không đồng nhất, đặc biệt là các thiết bị biên (edge devices) và nút mạng.

Tuy nhiên, hầu hết các khung giải pháp FD hiện tại vẫn phụ thuộc vào một máy chủ tổng hợp tập trung. Trong các triển khai quy mô lớn, máy chủ này trở thành một nút thắt về tin cậy, một điểm lỗi duy nhất và mục tiêu ưa thích của kẻ tấn công. Học liên kết phân tán (Decentralized Federated Learning -- DFL) giải quyết hạn chế này bằng cách phân phối trách nhiệm điều phối giữa các clients tham gia và loại bỏ sự phụ thuộc vào máy chủ vĩnh viễn \cite{braintorrent}.

Ngoài các vấn đề về truyền thông và điều phối, các hệ thống liên kết vẫn dễ bị tổn thương trước các tác nhân Byzantine cố tình đầu độc dữ liệu cục bộ, bản cập nhật mô hình hoặc các dự đoán được trao đổi \cite{tolpegin2020targetedflip, fang2020attacks, roux2025byzantine}. Các cuộc tấn công gần đây đã tiến hóa từ việc tạo ra các giá trị ngoại lai rõ ràng đến việc tạo ra các bản cập nhật độc hại gần giống với hành vi bất thường lành tính, khiến việc phát hiện ngày càng khó khăn \cite{roux2025byzantine}. Các phương pháp chống chịu Byzantine hiện có hầu hết tập trung vào FL trong không gian tham số, trong khi các phòng thủ mạnh mẽ cho FD phân tán vẫn còn chưa được khám phá đầy đủ.

Để giải quyết các hạn chế này, DFD-IDS được đề xuất như một khung giải pháp chắt lọc kiến thức liên kết phân tán có khả năng chống chịu đầu độc, với cơ chế kết hợp logits dựa trên năng lượng. DFD-IDS kết hợp Chắt lọc kiến thức dựa trên bằng chứng (Evidential Knowledge Distillation -- EKD) để bảo toàn độ bất định dự đoán, Cơ chế bỏ phiếu trắng dựa trên năng lượng (Energy-based Vote Abstention -- EVA) để lọc các mẫu công khai không đáng tin cậy, và Bộ lọc Blom Robust mới có khả năng chống chịu Byzantine. Thêm vào đó, chiến lược chọn điều phối viên phân tán được thực hiện dựa trên điểm số ngoại lai của từng client. Bằng cách kết hợp các thành phần này, khung giải pháp đề xuất cho phép chia sẻ tri thức hiệu quả giữa các clients không đồng nhất đồng thời duy trì khả năng chống chịu trước các cuộc tấn công đầu độc trong môi trường phân tán.

Các đóng góp chính của đề tài bao gồm:
\begin{itemize}
    \item Đề xuất DFD-IDS, một khung giải pháp chắt lọc kiến thức liên kết phi tập trung cho phát hiện xâm nhập cộng tác mà không cần máy chủ trung tâm vĩnh viễn.
    \item Giới thiệu cơ chế Energy-based Vote Abstention (EVA) giúp thích ứng phát hiện ngoài phân phối dựa trên năng lượng vào việc chọn lọc mẫu công khai trong môi trường liên kết.
    \item Phát triển Bộ lọc Blom Robust Filter giúp kết hợp logits chống chịu Byzantine trong môi trường clients không đồng nhất.
    \item Đánh giá DFD-IDS trước các cuộc tấn công đầu độc gần đây và chứng minh khả năng chống chịu vượt trội so với các phương pháp cơ sở hiện có.
\end{itemize}

\section{Mục tiêu, đối tượng và phạm vi}
\label{sec:objectives}

\subsection{Mục tiêu nghiên cứu}
Mục tiêu chính của đề tài là nghiên cứu, thiết kế và hiện thực một hệ thống phát hiện xâm nhập mạng dựa trên học liên kết phân tán, có khả năng chống chịu trước các cuộc tấn công đầu độc (poisoning attacks), hỗ trợ kiến trúc mô hình không đồng nhất (heterogeneous architectures) và giảm thiểu chi phí truyền thông. Cụ thể:

\begin{itemize}
    \item Nghiên cứu tổng quan về các hệ thống IDS dựa trên học máy và học liên kết.
    \item Tìm hiểu các phương pháp Federated Distillation và khả năng ứng dụng trong bảo mật mạng.
    \item Phân tích các dạng tấn công độc hại trong học liên kết và cơ chế phòng thủ tương ứng.
    \item Đề xuất và hiện thực khung giải pháp DFD-IDS với các thành phần chống chịu Byzantine.
    \item Đánh giá hiệu quả của giải pháp thông qua các kịch bản thực nghiệm.
\end{itemize}

\subsection{Đối tượng nghiên cứu}
Đối tượng nghiên cứu của đề tài bao gồm:
\begin{itemize}
    \item Hệ thống phát hiện xâm nhập mạng (Network-based IDS).
    \item Học liên kết tập trung và phân tán (Centralized \& Decentralized Federated Learning).
    \item Các phương pháp chắt lọc kiến thức liên kết (Federated Distillation).
    \item Các kỹ thuật phát hiện và chống chịu tấn công độc hại trong môi trường liên kết.
    \item Mô hình học sâu (Deep Neural Networks, GRU, CNN) cho dữ liệu mạng.
\end{itemize}

\subsection{Phạm vi nghiên cứu}
Đề tài tập trung vào:
\begin{itemize}
    \item Dữ liệu mạng từ các bộ dữ liệu chuẩn như CSE-CIC-IDS2018 và các tập dữ liệu liên quan.
    \item Môi trường thực nghiệm trong phòng lab với các máy chủ được cấu hình phù hợp.
    \item Đánh giá trên các chỉ số: Accuracy, Precision, Recall, F1-Score, chi phí truyền thông và khả năng chống chịu tấn công.
    \item Không bao gồm triển khai trên môi trường sản xuất thực tế.
\end{itemize}

\section{Những điểm mới của đề tài}
\label{sec:novelty}

Đề tài nghiên cứu có những điểm mới sau đây so với các công trình hiện có:

\begin{enumerate}
    \item \textbf{Kiến trúc phân tán hoàn toàn}: Không giống như các phương pháp FL-IDS hiện tại chủ yếu dựa trên server trung tâm, DFD-IDS loại bỏ hoàn toàn sự phụ thuộc vào permanent server thông qua cơ chế Dynamic Server Selection. Điều này giúp loại bỏ single point of failure và tăng tính sẵn sàng của hệ thống.
    \item \textbf{Energy-based Vote Abstention (EVA)}: Lần đầu tiên áp dụng cơ chế phát hiện out-of-distribution dựa trên năng lượng (Energy-based OOD) vào bài toán chọn lọc public samples trong Federated Distillation. EVA nhẹ hơn so với phương pháp discriminator-based của SSFL-IDS và không yêu cầu huấn luyện mô hình phụ trợ.
    \item \textbf{Blom Robust Filter}: Đề xuất một phương pháp robust aggregation mới dựa trên Blom order statistics, khắc phục hạn chế của phương pháp ngưỡng cố định (fixed threshold) trong Cronus. Phương pháp này có khả năng thích ứng với số lượng clients và số chiều logits khác nhau.
    \item \textbf{Kết hợp EKD và FL}: Ứng dụng Evidential Knowledge Distillation trong bối cảnh Federated Distillation để bảo toàn thông tin bất định (uncertainty) trong quá trình truyền thụ kiến thức, vượt trội hơn so với softmax-based distillation truyền thống.
\end{enumerate}

\section{Phương pháp nghiên cứu}
\label{sec:methodology}

Phương pháp nghiên cứu được áp dụng trong đề tài bao gồm:
\begin{itemize}
    \item \textbf{Nghiên cứu lý thuyết}: Thu thập, phân tích và tổng hợp các tài liệu khoa học liên quan đến IDS, Federated Learning, Federated Distillation, Byzantine-robust aggregation và các lĩnh vực liên quan.
    \item \text{Nghiên cứu thực nghiệm}: Thiết kế và hiện thực các mô hình, tiến hành thực nghiệm trên các bộ dữ liệu chuẩn, so sánh kết quả với các phương pháp nền tảng.
    \item \textbf{Đánh giá định lượng}: Sử dụng các chỉ số đánh giá khách quan để đo lường hiệu suất của hệ thống.
    \item \textbf{Phân tích so sánh}: So sánh phương pháp đề xuất với các phương pháp hiện có (FedAvg, SSFL-IDS, Cronus, FLAME, FLTrust) trên cùng điều kiện thực nghiệm.
\end{itemize}

\section{Kế hoạch thực tập chi tiết}
\label{sec:plan}

Quá trình thực tập tại Trung tâm An ninh mạng được chia thành 8 tuần với các mục tiêu cụ thể cho từng giai đoạn. Bảng~\ref{tab:plan_summary} tóm tắt kế hoạch thực tập tổng thể.

\begin{table}[H]
\centering
\caption{Kế hoạch thực tập tổng thể}
\label{tab:plan_summary}
\begin{tabular}{|c|c|c|}
\hline
\textbf{Giai đoạn} & \textbf{Tuần} & \textbf{Mục tiêu} \\
\hline
Định hướng & 1-2 & Làm quen môi trường, nghiên cứu tổng quan \\
\hline
Nghiên cứu lý thuyết & 3-4 & Tìm hiểu FL, FD, Byzantine attacks, defenses \\
\hline
Phát triển giải pháp & 5-6 & Thiết kế và hiện thực DFD-IDS \\
\hline
Thực nghiệm \& Đánh giá & 7 & Tiến hành thực nghiệm, thu thập kết quả \\
\hline
Hoàn thiện & 8 & Phân tích kết quả, viết báo cáo \\
\hline
\end{tabular}
\end{table}

\section{Cấu trúc báo cáo}
\label{sec:structure}

Báo cáo thực tập được tổ chức thành 5 chương:

\begin{itemize}
    \item \textbf{Chương 1 -- Tổng quan}: Giới thiệu về Trung tâm An ninh mạng, bối cảnh nghiên cứu, lý do chọn đề tài, mục tiêu, đối tượng, phạm vi và phương pháp nghiên cứu.
    \item \textbf{Chương 2 -- Cơ sở lý thuyết}: Trình bày các kiến thức nền tảng bao gồm tấn công APT, hệ thống IDS, học liên kết, chắt lọc kiến thức liên kết, tấn công độc hại, phòng thủ Byzantine, học máy khả giải và mạng SDN.
    \item \textbf{Chương 3 -- Phương pháp thực hiện}: Mô tả chi tiết khung giải pháp DFD-IDS được đề xuất, bao gồm kiến trúc tổng thể, cơ chế Energy-based Vote Abstention, Blom Robust Filter, cơ chế chọn server phân tán và Evidential Knowledge Distillation.
    \item \textbf{Chương 4 -- Thực nghiệm và đánh giá}: Trình bày môi trường thực nghiệm, tập dữ liệu, các kịch bản thử nghiệm, kết quả đạt được và bảng so sánh với các phương pháp khác.
    \item \textbf{Chương 5 -- Kết luận và hướng phát triển}: Tổng kết các kết quả đạt được, đánh giá ưu nhược điểm của giải pháp và đề xuất các hướng phát triển trong tương lai.
\end{itemize}

\pagebreak

%% ============================================================
%%  CHƯƠNG 2: CƠ SỞ LÝ THUYẾT
%% ============================================================
\chapter{Cơ sở lý thuyết}

\section{Tấn công APT}
\label{sec:apt}

\subsection{Tổng quan}
Advanced Persistent Threat (APT) là một dạng tấn công mạng tinh vi, kéo dài và có chủ đích, thường do các nhóm tác nhân có tổ chức thực hiện với mục tiêu xâm nhập và duy trì sự hiện diện lâu dài trong hệ thống của nạn nhân. Khác với các cuộc tấn công thông thường, APT thường nhắm vào các tổ chức lớn, cơ quan chính phủ, doanh nghiệp trong lĩnh vực tài chính, năng lượng và quốc phòng.

Đặc điểm chính của tấn công APT bao gồm:
\begin{itemize}
    \item Tính bền bỉ (Persistence): Kẻ tấn công duy trì truy cập trái phép trong thời gian dài, có thể kéo dài hàng tháng hoặc hàng năm.
    \item Tính tinh vi (Sophistication): Sử dụng nhiều kỹ thuật tiên tiến, khả năng tùy biến và thích ứng cao để vượt qua các lớp bảo vệ.
    \item Tính có chủ đích (Targeted): Tấn công nhắm vào một tổ chức hoặc cá nhân cụ thể, không phải tấn công hàng loạt.
    \item Tính ẩn danh (Stealth): Hoạt động âm thầm, tránh bị phát hiện bởi các hệ thống giám sát thông thường.
\end{itemize}

\subsection{Vòng đời tấn công APT}
Vòng đời của một cuộc tấn công APT thường trải qua nhiều giai đoạn:
\begin{enumerate}
    \item \textbf{Reconnaissance}: Thu thập thông tin về mục tiêu, bao gồm cấu trúc mạng, hệ thống, nhân sự và quy trình hoạt động.
    \item \textbf{Initial Compromise}: Xâm nhập ban đầu thông qua các kỹ thuật như spear-phishing, khai thác lỗ hổng zero-day hoặc tấn công watering hole.
    \item \textbf{Establish Foothold}: Thiết lập backdoor, cài đặt phần mềm độc hại và duy trì truy cập từ xa.
    \item \textbf{Escalate Privileges}: Leo thang đặc quyền để đạt được quyền truy cập cao nhất trong hệ thống.
    \item \textbf{Lateral Movement}: Di chuyển ngang trong mạng nội bộ để mở rộng phạm vi kiểm soát.
    \item \textbf{Maintain Persistence}: Duy trì sự hiện diện lâu dài thông qua nhiều cơ chế backdoor và kỹ thuật che giấu.
    \item \textbf{Exfiltration}: Thu thập và đánh cắp dữ liệu nhạy cảm.
\end{enumerate}

\subsection{Cách phòng chống}
Phòng chống tấn công APT đòi hỏi một chiến lược đa lớp (defense-in-depth), kết hợp nhiều giải pháp như:
\begin{itemize}
    \item Hệ thống phát hiện xâm nhập (IDS/IPS) kết hợp học máy.
    \item Phân tích hành vi mạng và thiết bị (UEBA).
    \item Săn tìm mối đe dọa chủ động (Threat Hunting).
    \item Chia sẻ thông tin tình báo mối đe dọa (CTI).
    \item Kiểm soát truy cập và quản lý đặc quyền (PAM).
    \item Đào tạo nhận thức an ninh cho nhân viên.
\end{itemize}

\section{Hệ thống IDS}
\label{sec:ids}

Hệ thống phát hiện xâm nhập (IDS) là một thành phần quan trọng trong hạ tầng an ninh mạng hiện đại, cho phép xác định các hoạt động độc hại trong các mạng doanh nghiệp, đám mây, công nghiệp và IoT. Như đã trình bày trong Mục~\ref{sec:reason}, bối cảnh mối đe dọa tiếp tục phát triển với các cuộc tấn công ngày càng tinh vi.

\subsection{Phân loại IDS}
IDS có thể được phân loại theo nhiều tiêu chí khác nhau:

\textbf{Theo phương pháp phát hiện:}
\begin{itemize}
    \item \textbf{Signature-based IDS}: Phát hiện dựa trên các mẫu tấn công đã biết (chữ ký). Hiệu quả cao với các tấn công cũ nhưng không phát hiện được tấn công mới (zero-day).
    \item \textbf{Anomaly-based IDS}: Phát hiện dựa trên sự khác biệt so với hành vi bình thường của hệ thống. Có thể phát hiện tấn công mới nhưng tỷ lệ dương tính giả cao.
\end{itemize}

\textbf{Theo nguồn dữ liệu:}
\begin{itemize}
    \item \textbf{Network-based IDS (NIDS)}: Phân tích lưu lượng mạng để phát hiện tấn công.
    \item \textbf{Host-based IDS (HIDS)}: Giám sát hoạt động trên một máy chủ cụ thể, bao gồm log hệ thống, tiến trình và tệp tin.
\end{itemize}

\subsection{NIDS dựa trên học sâu}
Những tiến bộ gần đây trong ML và DL đã cải thiện đáng kể khả năng phát hiện của IDS thông qua việc tự động học các mẫu lưu lượng phức tạp từ dữ liệu mạng quy mô lớn. Trong số nhiều kiến trúc học sâu, Mạng nơ-ron tích chập (CNN) có hiệu quả trong việc trích xuất các đặc trưng không gian từ luồng mạng, trong khi Bộ nhớ hồi tiếp có cổng (GRU) vượt trội trong việc nắm bắt các phụ thuộc thời gian. Mạng nơ-ron đồ thị (GNN) cũng đã nổi lên như một phương pháp tiếp cận đầy hứa hẹn để mô hình hóa cấu trúc liên kết mạng và các mối quan hệ luồng dữ liệu.

\subsubsection{Mạng CNN trong phát hiện xâm nhập}
Convolutional Neural Networks (CNN) có khả năng tự động học các đặc trưng không gian từ dữ liệu thô thông qua các lớp tích chập (convolutional layers). Trong bài toán IDS, lưu lượng mạng có thể được biểu diễn dưới dạng ma trận 2 chiều (ví dụ: time $\times$ features) và CNN có thể học các mẫu tấn công cục bộ trong không gian đặc trưng này. Các nghiên cứu gần đây cho thấy CNN đạt hiệu quả cao trong việc phân loại lưu lượng mạng và phát hiện các cuộc tấn công DDoS, botnet và mã độc.

\subsubsection{Mạng GRU trong phát hiện xâm nhập}
Gated Recurrent Units (GRU) là một biến thể của mạng nơ-ron hồi tiếp (RNN), được thiết kế để giải quyết vấn đề gradient biến mất (vanishing gradient) thông qua cơ chế cổng (gate mechanism). Trong bối cảnh IDS, GRU có khả năng mô hình hóa các phụ thuộc thời gian trong chuỗi các gói tin mạng, cho phép phát hiện các hành vi tấn công có tính chất tuần tự như port scanning, brute force và các giai đoạn của tấn công APT.

\subsubsection{Mạng GNN trong phát hiện xâm nhập}
Graph Neural Networks (GNN) mở rộng học sâu lên dữ liệu có cấu trúc đồ thị. Trong mạng máy tính, các thiết bị và luồng dữ liệu có thể được mô hình hóa như một đồ thị, nơi các đỉnh là các thiết bị mạng và các cạnh là các kết nối hoặc luồng dữ liệu giữa chúng. GNN có thể học các biểu diễn của các nút và cạnh trong đồ thị này, phát hiện các mẫu tấn công phức tạp và tấn công APT có tính chất lan truyền ngang.

\subsubsection{Thách thức trong triển khai}
Thách thức chính trong việc triển khai IDS dựa trên ML là yêu cầu về lượng lớn dữ liệu huấn luyện được gán nhãn, điều thường khó đạt được do các ràng buộc về quyền riêng tư và quy định. Ngoài ra, các mẫu lưu lượng mạng thay đổi theo thời gian, đòi hỏi các bản cập nhật mô hình liên tục. Chi phí tính toán để huấn luyện các mô hình học sâu cũng là một vấn đề đối với các môi trường hạn chế tài nguyên. Những thách thức này thúc đẩy việc áp dụng Học liên kết (Federated Learning), cho phép huấn luyện mô hình cộng tác trong khi vẫn bảo vệ quyền riêng tư dữ liệu, như sẽ được thảo luận trong các phần tiếp theo.

\section{Học liên kết (Federated Learning)}
\label{sec:fl}

\subsection{Học liên kết tập trung (Centralized Federated Learning)}
Trong Học liên kết tập trung (Centralized FL), một máy chủ khởi tạo và duy trì một mô hình toàn cục trong khi các clients phân tán cùng tối ưu hóa mô hình này bằng dữ liệu riêng của họ. Khung giải pháp được công nhận rộng rãi nhất trong bối cảnh này là FedAvg \cite{fedavg}. Trong FedAvg, mô hình toàn cục $f^0$ được khởi tạo bởi máy chủ và gửi đến tất cả clients trong tập $K$. Với mỗi vòng liên kết $t$, các clients đồng thời tối ưu hóa các tham số $w^t$ của $f^t$ sử dụng tập dữ liệu riêng $D_k, k\in K$, $w_k^t \leftarrow ClientUpdate(k, w_t)$ trước khi truyền các tham số mô hình cục bộ $w_k^t$ đến máy chủ. Máy chủ sau đó thực hiện tổng hợp có trọng số trên tập các tham số client để tính toán tham số mô hình toàn cục cho vòng tiếp theo $w^{t+1}$ theo Phương trình~\ref{fedavg}.

\begin{equation}
\label{fedavg}
w^{t+1} = \sum_{k \in K}\frac{|D_k|}{|D|}w_k^t
\end{equation}

trong đó $D_k$ là tập dữ liệu cục bộ của client $k$, $|D_k|$ là kích thước của nó, và $|D|$ là tổng kích thước mẫu của tất cả clients tham gia trong vòng liên kết hiện tại $t$. Máy chủ gửi mô hình toàn cục đã tính toán đến tất cả clients và lặp lại sang vòng tiếp theo.

Trong bối cảnh phát hiện xâm nhập, mỗi client quan sát các khối lượng mẫu tấn công khác nhau; do đó, sự phân phối dữ liệu không đồng nhất cao cả về số lượng và loại hình là một thách thức cố hữu đối với học liên kết. Điều kiện này, được biết đến rộng rãi với tên gọi dữ liệu không độc lập và đồng nhất (Non-IID), gây ra hiện tượng ``client drift'' \cite{he2022fedssd}. Với dữ liệu Non-IID, mục tiêu tối ưu hóa cục bộ của mỗi client không thể tương thích với mục tiêu toàn cục do sự thiên lệch cục bộ của tập dữ liệu riêng \cite{acar2021feddyn}, thường làm chậm hội tụ và giảm hiệu suất mô hình toàn cục.

Để giải quyết thách thức Non-IID, nhiều thuật toán FL đã được đề xuất. Các phương pháp được công nhận như FedProx \cite{li2020fedprox}, FedMLB \cite{kim2022fedmlb} và FedDyn \cite{acar2021feddyn} điều chỉnh mục tiêu tối ưu hóa cục bộ để kiềm chế client drift, cải thiện khả năng chống chịu với dữ liệu Non-IID.

Mặc dù có khả năng chống chịu cao với điều kiện Non-IID, FL tập trung yêu cầu các client truyền tham số mô hình, gây ra những hạn chế nghiêm trọng trong triển khai thực tế. Các nghiên cứu gần đây như Deep Leakage from Gradients (DLG) \cite{zhu2019dlg} và improved Deep Leakage from Gradients (iDLG) \cite{zhao2020idlg} đã thành công trong việc tái tạo mẫu huấn luyện từ gradient với độ chính xác đáng báo động, trong khi \cite{geiping2020inverting} đề xuất một cuộc tấn công tái tạo mạnh hơn thích ứng trực tiếp với môi trường FL, cho thấy rằng các tham số mô hình không tự nhiên bảo vệ quyền riêng tư trước một máy chủ độc hại. Ngoài tấn công đảo ngược mô hình, nhiều công trình đề xuất tấn công suy diễn \cite{pyrgelis2017knockknock, melis2019mia, ganju2018inference} trên các tham số mô hình.

Nhiều công trình đã tìm cách giải quyết điểm yếu này: SecureAggregation \cite{bonawitz2016secagg} và DeepFed \cite{li2020deepfed} là các biến thể bảo vệ quyền riêng tư dựa trên cơ chế mật mã như Diffie-Hellman masking hoặc Pallier Homomorphic Encryption. Tuy nhiên, chúng tạo ra chi phí mật mã và các vấn đề thực tế đáng kể. Ví dụ, SecureAggregation tạo thêm các vòng truyền thông và cơ chế phục hồi dropout, làm tăng độ phức tạp giao thức trong triển khai quy mô lớn. Tương tự, các phương pháp dựa trên mã hóa đồng hình như DeepFed tạo ra chi phí tính toán và quản lý khóa đáng kể.

\subsection{Chắt lọc kiến thức liên kết tập trung (Centralized Federated Distillation)}
\label{sec:fd}
Chắt lọc kiến thức liên kết (Federated Distillation -- FD) \cite{li2019fedmd, tan2022fedproto, jeong2018federateddistillation} nổi lên như một biến thể của Học liên kết kết hợp FL với Chắt lọc kiến thức (Knowledge Distillation) \cite{hinton2015kd}. Các khung giải pháp FD, bao gồm FedMD \cite{li2019fedmd} và FederatedDistillation \cite{jeong2018federateddistillation} truyền thông các biểu diễn tri thức trung gian như logits hoặc class prototypes \cite{tan2022fedproto} thay vì tham số mô hình.

FederatedDistillation tính toán vector logit trung bình đã kích hoạt softmax $\bar{p}_{c,k}$ của mỗi lớp $c$ trong tập dữ liệu riêng $D_k$ của mỗi client $k$, $\bar{p}_{c,k} \leftarrow \frac{1}{|c|}\sum_{x \in D_k, y|x = c}p(x)$. Sau đó, máy chủ tính tổng các vector logit của tất cả clients cho mỗi lớp $c \in C$, $\bar{p}_c \leftarrow \sum_{k \in K}\bar{p}_{c,k}$. Với mỗi lớp, mỗi client trừ vector trung bình của chính nó khỏi vector toàn cục để thu được vector trung bình leave-one-out $\hat{p}_{c,k} \leftarrow \frac{1}{|K|-1} (\bar{p}_c - \bar{p}_{c,k})$. Mỗi client tối ưu hóa mô hình cục bộ qua Phương trình~\ref{federateddistillation} cho mỗi mẫu $x$ trong tập dữ liệu riêng:

\begin{equation}
\label{federateddistillation}
w_k \leftarrow w_k - \eta \nabla\{\phi(p(x),y_i) + \gamma \cdot \phi(p(x), \hat{p}_{c,k}\}
\end{equation}

trong đó $\phi$ là hàm mất mát cross-entropy, $\gamma$ là siêu tham số điều khiển cường độ chắt lọc.

Các khung giải pháp FD gần đây đã phát triển để hoạt động trên một tập dữ liệu công khai $D_p$, có thể truy cập bởi tất cả các thành viên liên kết. FedMD \cite{li2019fedmd} là một công trình phổ biến sử dụng tập dữ liệu công khai có nhãn. Trong cách tiếp cận của họ, tất cả clients huấn luyện mô hình riêng trên tập dữ liệu công khai cho đến khi hội tụ, sau đó chuyển sang tập dữ liệu riêng trước quá trình liên kết. Với mỗi vòng liên kết, clients tính toán raw logits cho mỗi mẫu của tập dữ liệu công khai. Máy chủ tổng hợp các logits của clients để tạo thành một consensus. Mỗi client huấn luyện mô hình cục bộ để tiến gần đến consensus này. Trong giai đoạn Digest, consensus của máy chủ đóng vai trò là mô hình teacher, trong khi mỗi mô hình client đóng vai trò là mô hình student. Mỗi mô hình student học cách khớp đầu ra của nó trên các mẫu công khai với mô hình teacher consensus thông qua một mục tiêu chắt lọc.

SSFL-IDS \cite{zhao2022ssfl} tận dụng các nguyên lý của FD và học bán giám sát (semi-supervised learning). SSFL-IDS không yêu cầu tập dữ liệu công khai phải được gán nhãn. Trong cách tiếp cận của họ, mỗi client tối ưu hóa mô hình cục bộ bằng tập dữ liệu riêng và tính toán logits đã kích hoạt softmax trên các mẫu công khai. SSFL-IDS triển khai một bộ phân loại nhị phân để loại bỏ các dự đoán trên các mẫu công khai mà mô hình chính không chắc chắn.

\subsection{Học liên kết phân tán (Decentralized Federated Learning)}
\label{sec:dfl}
Học liên kết phân tán (Decentralized Federated Learning -- DFL) loại bỏ sự phụ thuộc vào máy chủ trung tâm vĩnh viễn bằng cách phân phối trách nhiệm điều phối giữa các clients tham gia \cite{beltran2023dflsurvey}. Các khung giải pháp DFL hiện có được chia thành nhiều cấu trúc liên kết khác nhau, bao gồm kiến trúc vòng (ring), lưới (mesh), cụm (cluster) và hình sao (star). Các cấu trúc này khác nhau về độ phức tạp truyền thông, khả năng chịu dropout client, tốc độ hội tụ liên kết và chiến lược tổng hợp.

Trong số các cấu trúc này, DFL dạng hình sao duy trì mẫu truyền thông gần giống với FL tập trung nhất, nơi một điều phối viên tạm thời thực hiện tổng hợp trong mỗi vòng. Không giống như FL tập trung, vai trò điều phối viên được chỉ định động trong quá trình liên kết. BrainTorrent \cite{braintorrent} là một công trình tiêu biểu thể hiện một khung giải pháp FedAvg phân tán ngang hàng (peer-to-peer) theo cấu trúc hình sao.

Trong bối cảnh phát hiện xâm nhập, các nghiên cứu gần đây đã khám phá DFL-based IDS. 2DF-IDS \cite{friha20232dfl-ids} đề xuất một IDS dựa trên FL phân tán cho IIoT với cơ chế tổng hợp hoàn toàn phân tán. DOF-IDS \cite{nakip2023dfl-ids} cho phép FL phân tán cho các nút IoT và IP với cách tiếp cận dựa trên đăng ký. Tuy nhiên, mặc dù các khung giải pháp này loại bỏ máy chủ trung tâm vĩnh viễn, hầu hết vẫn tiếp tục trao đổi tham số mô hình giữa các thành viên, để lại chi phí truyền thông và rủi ro đảo ngược mô hình như những thách thức đáng kể.

\section{Tấn công độc hại (Poisoning Attacks) và Phòng thủ}
\label{sec:poisoning}

Bản chất cộng tác của học liên kết khiến quá trình huấn luyện dễ bị tấn công bởi các thành viên độc hại làm hỏng quá trình tổng hợp. Thông qua việc thao túng dữ liệu huấn luyện cục bộ, bản cập nhật mô hình và không gian dự đoán, các cuộc tấn công đầu độc có thể làm suy giảm đáng kể hiệu suất mô hình toàn cục và đặt ra thách thức lớn cho việc triển khai thực tế các hệ thống phát hiện xâm nhập liên kết.

\subsection{Tấn công đầu độc dữ liệu (Data Poisoning)}
Tấn công đầu độc dữ liệu là một cuộc tấn công có thể được thực hiện trong tất cả các khung giải pháp liên kết. Nó thao túng tập dữ liệu huấn luyện cục bộ của các client Byzantine trước khi huấn luyện, thúc đẩy quá trình tối ưu hóa cục bộ tạo ra một mô hình bị đầu độc. Một ví dụ phổ biến trong bối cảnh phát hiện xâm nhập là đảo nhãn có chủ đích (targeted label flipping), nơi các mẫu từ lớp tấn công thường xuyên xuất hiện được gán nhãn lại thành lưu lượng lành tính, buộc mô hình học các ranh giới quyết định sai. Các client bị đầu độc sau đó gửi các tham số mô hình bị đầu độc hoặc tạo ra các biểu diễn sai lệch ảnh hưởng xấu đến mô hình tổng hợp hoặc consensus.

\subsection{Tấn công đầu độc mô hình (Model Poisoning)}
Tấn công đầu độc mô hình sửa đổi các bản cập nhật mô hình cục bộ trước khi tổng hợp toàn cục, mà không cần bất kỳ dữ liệu huấn luyện nào. Kẻ tấn công có thể tạo ra các bản cập nhật tùy ý làm hỏng nghiêm trọng mô hình được tổng hợp một cách ngây thơ, khiến nó hoàn toàn không đáng tin cậy và làm đình trệ quá trình liên kết.

PoisonedFL \cite{xie2025poisonedfl} là một cuộc tấn công đầu độc mô hình gần đây có khả năng thích ứng với các phòng thủ hiện đại, bao gồm FLTrust \cite{cao2020fltrust} và FLAME \cite{nguyen2022flame}. Tất cả các client Byzantine trong cuộc tấn công này tạo ra một bản cập nhật bị đầu độc giống hệt nhau để giảm hiệu ứng ``tự triệt tiêu'' (self cancellation), đồng thời thích ứng bằng cách giảm cường độ tấn công nếu phát hiện sự hiện diện của cơ chế phòng thủ.

\subsection{Tấn công đầu độc Logit (Logit Poisoning)}
Trong các hệ thống chắt lọc kiến thức liên kết, các client Byzantine trực tiếp thao túng các bản cập nhật biểu diễn ảnh hưởng xấu đến consensus tổng hợp. Tấn công tối đa hóa mất mát (Loss Maximization Attack -- LMA) được giới thiệu bởi \cite{roux2025byzantine} tạo ra các logits đối nghịch nhằm tối đa hóa sự bất đồng trong quá trình tổng hợp, làm hỏng mô hình teacher consensus, và từ đó làm hỏng các mô hình student vốn là các mô hình lành tính.

Các cuộc tấn công đầu độc gần đây đã tiến hóa từ việc làm suy giảm FL bằng các giá trị ngoại lai cực đoan rõ ràng đến việc tạo ra các bản cập nhật độc hại ẩn danh nằm gần với hành vi Non-IID lành tính. Kết quả là, việc phân biệt các thành viên độc hại khỏi các client không đồng nhất tự nhiên ngày càng trở nên khó khăn.

\subsection{Cơ chế phòng thủ (Anti-Poisoning)}
\label{sec:defense}

Để giảm thiểu rủi ro từ các cuộc tấn công Byzantine, nhiều công trình nghiên cứu FL đã đề xuất các cơ chế khác nhau nhằm giảm hoặc loại bỏ hoàn toàn ảnh hưởng độc hại của các bản cập nhật trong quá trình tổng hợp.

\textbf{Các Robust Aggregation Rule cổ điển:}
\begin{itemize}
    \item \textbf{Coordinate-wise Median} \cite{yin2018fedcomed}: Chọn giá trị trung vị cho mỗi tham số, loại bỏ các giá trị ngoại lai cực đoan.
    \item \textbf{TrimmedMean} \cite{yin2018fedcomed}: Loại bỏ phần cao nhất và thấp nhất của các bản cập nhật client cho mỗi tham số.
    \item \textbf{Krum} \cite{blanchard2017krum}: Chọn vector tham số gần nhất với tất cả các lân cận của nó.
    \item \textbf{Bulyan} \cite{guerraoui2018bulyan}: Kết hợp Multi-Krum và TrimmedMean trong một quy trình tổng hợp hai giai đoạn.
\end{itemize}

\textbf{Các phương pháp hiện đại:}
\begin{itemize}
    \item \textbf{FLAME} \cite{nguyen2022flame}: Máy chủ tính khoảng cách cosine từng cặp cho mỗi bản cập nhật client, phân cụm chúng bằng HDBSCAN \cite{mcinnes2017hdbscan}, loại bỏ các cụm ngoại lai, cắt tỉa các mô hình được chấp nhận và áp dụng nhiễu thích ứng.
    \item \textbf{FLTrust} \cite{cao2020fltrust}: Giả định một tập dữ liệu công khai có nhãn là tập dữ liệu gốc. Mô hình toàn cục được huấn luyện trên tập dữ liệu gốc này song song với huấn luyện cục bộ của client, và độ tương đồng cosine được cắt bằng ReLU được sử dụng để tính điểm tin cậy.
    \item \textbf{Cronus} \cite{chang2019cronus}: Triển khai RobustFilter như một bộ tổng hợp logit chống chịu Byzantine, lọc các dự đoán ngoại lai trước khi gán nhãn giả (pseudo-labeling).
    \item \textbf{ExpGuard} \cite{roux2025byzantine}: Đo khoảng cách của mỗi bản cập nhật client so với kết quả tổng hợp của các luật tổng hợp mạnh mẽ để tính điểm ngoại lai $\sigma_k$ cho mỗi client.
\end{itemize}

\subsection{Chắt lọc kiến thức dựa trên bằng chứng (Evidential Knowledge Distillation)}
\label{sec:ekd}

Chắt lọc kiến thức truyền thống truyền thụ tri thức thông qua các phân phối xác suất thu được bằng chuẩn hóa softmax. Mặc dù hiệu quả, softmax loại bỏ thông tin chứa trong độ lớn tuyệt đối của các logits, khiến các vector dự đoán với mức độ tin cậy khác nhau tạo ra các phân phối xác suất giống hệt nhau sau chuẩn hóa.

For example, since softmax satisfies $\text{softmax}(z+c) = \text{softmax}(z)$ for all $c \in \mathbb{R}$, assume we have two logit vectors of different magnitudes:

$$z^{(1)} = [0.1, -0.1, 0]$$
$$z^{(2)} = z^{(1)}+4 = [4.1, 3.9, 4]$$

$\text{softmax}(z^{(1)}) = \text{softmax}(z^{(2)}) \approx [0.367, 0.301, 0.332]$. Chuẩn hóa softmax bất biến với các phép dịch chuyển cộng, do đó thông tin về độ lớn tuyệt đối của các logits bị loại bỏ.

Chắt lọc kiến thức dựa trên bằng chứng (Evidential Knowledge Distillation -- EKD), được đề xuất bởi \cite{xiang2025ekd}, truyền độ bất định của mô hình teacher cùng với kiến thức của nó đến các mô hình student. Trong EKD, các raw logits được kích hoạt bằng hàm mũ $e=\text{exp}(z)$, kết hợp với trọng số tiên nghiệm $\lambda$ để tạo ra vector tham số Dirichlet $|C|$ chiều $\alpha = e + \lambda$. Mô hình student chắt lọc kiến thức từ mô hình teacher bằng cách tối ưu hóa:

\begin{equation}
\label{ekd}
\begin{split}
L & = L_{1st} + \gamma L_{2nd} \\
  & = KL\left(\frac{\alpha_i^T}{\sum_{i \in C}\alpha_i^T} \middle|\middle| \frac{\alpha_i^S}{\sum_{i \in C}\alpha_i^S}\right) \\
  & +\gamma \cdot KL(Dir(p^T|\alpha^T)||Dir(p^S|\alpha^S))
\end{split}
\end{equation}

\subsection{Phát hiện ngoài phân phối dựa trên năng lượng (Energy-based OOD Detection)}
\label{sec:energy}

Ước lượng độ tin cậy là rất quan trọng trong việc xác định các dự đoán không chắc chắn. Các phương pháp truyền thống dựa trên xác suất softmax tối đa tạo ra các dự đoán quá tự tin ngay cả đối với các mẫu ngoài phân phối \cite{liu2020energy}. Phương pháp phát hiện OOD dựa trên năng lượng được đề xuất bởi \cite{liu2020energy} định lượng độ bất định của mô hình trên một mẫu dữ liệu $x$ dưới dạng năng lượng:

\begin{equation}
\label{energy}
E(x) = -T \log\sum_i e^{{z_i}/T}
\end{equation}

Năng lượng thấp cho thấy dự đoán tin cậy, trong khi giá trị năng lượng cao thường cho thấy một mẫu không chắc chắn hoặc ngoài phân phối.

\section{Học máy khả giải (Explainable AI)}
\label{sec:xai}

Học máy khả giải (Explainable Artificial Intelligence -- XAI) là lĩnh vực nghiên cứu nhằm làm cho các mô hình học máy trở nên minh bạch và có thể giải thích được. Trong bối cảnh an ninh mạng, khả năng diễn giải các quyết định của hệ thống IDS là vô cùng quan trọng, giúp các nhà phân tích bảo mật hiểu được lý do tại sao một cảnh báo được đưa ra và có thể phản ứng phù hợp.

Tầm quan trọng của XAI trong bảo mật mạng thể hiện ở một số khía cạnh sau:
\begin{itemize}
    \item \textbf{Tăng độ tin cậy}: Giải thích được các dự đoán của mô hình giúp nhà phân tích tin tưởng hơn vào hệ thống IDS, đặc biệt là khi cần đưa ra quyết định ngăn chặn tấn công tự động.
    \item \text{Hỗ trợ điều tra sự cố}: Khi một cảnh báo được kích hoạt, việc biết được đặc trưng nào đã dẫn đến quyết định đó giúp rút ngắn thời gian điều tra và khắc phục.
    \item \textbf{Cải thiện mô hình}: Phân tích các giải thích có thể phát hiện ra các lỗi tiềm ẩn của mô hình, chẳng hạn như học các đặc trưng nhiễu hoặc bị tấn công đầu độc.
    \item \textbf{Tuân thủ quy định}: Nhiều quy định về bảo vệ dữ liệu và an ninh mạng yêu cầu khả năng giải thích các quyết định tự động.
\end{itemize}

\subsection{Phân loại}
Các phương pháp XAI có thể được phân loại theo nhiều tiêu chí:

\textbf{Theo phạm vi giải thích:}
\begin{itemize}
    \item \textbf{Global explanation}: Giải thích hành vi tổng thể của mô hình.
    \item \textbf{Local explanation}: Giải thích cho một dự đoán cụ thể.
\end{itemize}

\textbf{Theo phương pháp:}
\begin{itemize}
    \item \textbf{Intrinsic methods}: Mô hình tự thân có khả năng giải thích (ví dụ: decision tree, linear regression).
    \item \textbf{Post-hoc methods}: Giải thích sau khi mô hình đã được huấn luyện (ví dụ: LIME, SHAP, GNNExplainer).
\end{itemize}

\subsection{Các phương pháp phổ biến}
\begin{itemize}
    \item \textbf{SHAP (SHapley Additive exPlanations)}: Dựa trên lý thuyết trò chơi, tính toán giá trị Shapley để xác định mức độ đóng góp của mỗi đặc trưng vào dự đoán. SHAP cung cấp cả giải thích toàn cục và địa phương, có nền tảng lý thuyết vững chắc từ giá trị Shapley trong lý thuyết trò chơi hợp tác.
    \item \textbf{LIME (Local Interpretable Model-agnostic Explanations)}: Xấp xỉ mô hình phức tạp bằng một mô hình đơn giản (ví dụ: hồi quy tuyến tính) trong lân cận của mẫu cần giải thích. LIME hoạt động độc lập với kiến trúc mô hình, phù hợp với nhiều loại mô hình IDS khác nhau.
    \item \textbf{GNNExplainer}: Phương pháp giải thích dành riêng cho mô hình đồ thị (Graph Neural Networks), xác định các đỉnh và cạnh quan trọng nhất trong đồ thị. Đối với bài toán phát hiện xâm nhập trên đồ thị lưu lượng mạng, GNNExplainer có thể chỉ ra các luồng dữ liệu và kết nối đáng ngờ nhất.
\end{itemize}

\subsection{Ý nghĩa của XAI trong bảo mật}
Trong các hệ thống bảo mật, đặc biệt là hệ thống phát hiện xâm nhập, khả năng diễn giải dự đoán không chỉ là một tính năng bổ sung mà là một yêu cầu quan trọng. Khi một IDS phát hiện một cuộc tấn công, nhà phân tích bảo mật cần hiểu được cơ sở của cảnh báo đó để có thể:
\begin{itemize}
    \item Xác thực cảnh báo và loại trừ các trường hợp dương tính giả.
    \item Điều tra nguồn gốc và phương thức tấn công.
    \item Đưa ra biện pháp ứng phó phù hợp và kịp thời.
    \item Cải thiện các quy tắc phát hiện cho các lần tấn công tương tự trong tương lai.
\end{itemize}

Các phương pháp giải thích như SHAP và LIME có thể được áp dụng trực tiếp vào các mô hình IDS dựa trên học sâu, trong khi GNNExplainer phù hợp với các mô hình phân tích đồ thị nguồn gốc (provenance graphs) để phát hiện tấn công APT.

Việc áp dụng các phương pháp giải thích vào hệ thống IDS giúp tăng độ tin cậy, hỗ trợ điều tra sự cố và cải thiện khả năng phát hiện các mối đe dọa mới.

\section{Đồ thị nguồn gốc (Provenance Graph)}
\label{sec:provenance}

Đồ thị nguồn gốc (Provenance Graph) là một cấu trúc dữ liệu biểu diễn mối quan hệ nhân quả giữa các đối tượng trong hệ thống (tiến trình, tệp tin, socket, v.v.) và các hành động (đọc, ghi, thực thi, kết nối). Trong lĩnh vực an ninh mạng, provenance graph được sử dụng để theo dõi và phân tích chuỗi tấn công phức tạp, đặc biệt là tấn công APT.

\subsection{Tiềm năng trong phát hiện APT}
Provenance graph mang lại nhiều lợi thế trong việc phát hiện tấn công APT:
\begin{itemize}
    \item \textbf{Tính toàn vẹn về mặt nhân quả}: Provenance graph ghi lại đầy đủ mối quan hệ nhân quả giữa các sự kiện, giúp tái dựng toàn bộ chuỗi tấn công từ giai đoạn xâm nhập ban đầu đến giai đoạn đánh cắp dữ liệu.
    \item \textbf{Phát hiện hành vi bất thường}: Bằng cách phân tích cấu trúc của provenance graph, các hành vi bất thường như leo thang đặc quyền, di chuyển ngang và exfiltration có thể được phát hiện thông qua các bất thường về cấu trúc đồ thị.
    \item \textbf{Tích hợp với GNN}: Graph Neural Networks có thể được áp dụng trực tiếp trên provenance graph để học các biểu diễn và phát hiện các mẫu tấn công phức tạp.
    \item \textbf{Giảm thiểu dương tính giả}: Nhờ mối quan hệ nhân quả rõ ràng, provenance graph giúp loại bỏ các cảnh báo nhiễu bằng cách chỉ tập trung vào các chuỗi sự kiện có liên quan đến nhau.
\end{itemize}

\section{Mạng SDN}
\label{sec:sdn}

Software-Defined Networking (SDN) là một kiến trúc mạng cho phép tách biệt mặt phẳng điều khiển (control plane) và mặt phẳng dữ liệu (data plane), tập trung hóa khả năng quản lý và cấu hình mạng thông qua một bộ điều khiển trung tâm.

\subsection{Kiến trúc SDN}
Kiến trúc SDN bao gồm ba lớp chính:
\begin{itemize}
    \item \textbf{Lớp ứng dụng (Application Layer)}: Chứa các ứng dụng mạng như firewall, load balancer, IDS/IPS.
    \item \textbf{Lớp điều khiển (Control Layer)}: Bộ điều khiển SDN chịu trách nhiệm quản lý luồng dữ liệu và cấu hình thiết bị.
    \item \textbf{Lớp dữ liệu (Data Layer)}: Các thiết bị mạng vật lý hoặc ảo như switch, router thực hiện chuyển tiếp gói tin dựa trên chỉ thị từ lớp điều khiển.
\end{itemize}

\subsection{SDN trong phát hiện tấn công}
SDN mang lại tiềm năng lớn trong việc phát hiện và ngăn chặn tấn công mạng nói chung và tấn công APT nói riêng:
\begin{itemize}
    \item \textbf{Tập trung hóa giám sát}: Bộ điều khiển SDN có tầm nhìn toàn diện về toàn bộ lưu lượng mạng, giúp phát hiện các hành vi bất thường và tấn công đa giai đoạn một cách hiệu quả.
    \item \textbf{Linh hoạt trong cấu hình}: Có thể nhanh chóng thay đổi chính sách mạng để cô lập thiết bị bị tấn công, chặn lưu lượng độc hại hoặc chuyển hướng lưu lượng qua các thiết bị kiểm tra bảo mật.
    \item \textbf{Tích hợp với IDS}: Hệ thống IDS có thể tương tác trực tiếp với bộ điều khiển SDN thông qua các API để thực hiện các hành động ngăn chặn tự động (automatic threat mitigation), giảm thời gian phản ứng từ phút xuống còn mili-giây.
    \item \textbf{Phân tích lưu lượng toàn diện}: Với khả năng thu thập thông tin từ tất cả các switch trong mạng, bộ điều khiển SDN có thể xây dựng một bức tranh toàn cảnh về trạng thái mạng, hỗ trợ phát hiện các cuộc tấn công APT có tính chất lan truyền ngang.
\end{itemize}

\subsection{Tích hợp SDN với Federated Learning IDS}
Sự kết hợp giữa SDN và FL-based IDS mở ra những hướng tiếp cận mới cho an ninh mạng. Bộ điều khiển SDN có thể đóng vai trò là một federated client hoặc aggregation server trong kiến trúc FL. Các lợi ích chính bao gồm:
\begin{itemize}
    \item Tận dụng tầm nhìn toàn diện của SDN để cải thiện chất lượng phát hiện xâm nhập.
    \item Sử dụng khả năng lập trình của SDN để triển khai nhanh các chính sách bảo mật dựa trên kết quả phát hiện của IDS.
    \item Phân phối tải giám sát và phát hiện trên nhiều thiết bị mạng, giảm áp lực lên bộ điều khiển trung tâm.
    \item Tăng khả năng chống chịu của hệ thống thông qua kiến trúc phân tán.
\end{itemize}

\pagebreak

%% ============================================================
%%  CHƯƠNG 3: PHƯƠNG PHÁP THỰC HIỆN
%% ============================================================
\chapter{Phương pháp thực hiện}

Các khung giải pháp Chắt lọc kiến thức liên kết hiện tại hỗ trợ hiệu quả các kiến trúc mô hình không đồng nhất, nhưng chúng hoặc vẫn dễ bị thao túng Byzantine, phụ thuộc vào điều phối tập trung, hoặc thường xuyên đưa các mẫu công khai không chắc chắn vào quá trình chắt lọc. Để giải quyết các thách thức này, DFD-IDS được đề xuất như một khung giải pháp chắt lọc kiến thức liên kết phân tán kết hợp cơ chế bỏ phiếu trắng dựa trên năng lượng, Chắt lọc kiến thức dựa trên bằng chứng, kết hợp logits chống chịu đầu độc và bầu chọn máy chủ động. Trong DFD-IDS, các client tinh chỉnh mô hình cá nhân hóa và trao đổi các dự đoán được điều chỉnh bằng năng lượng trên một tập dữ liệu công khai chia sẻ. Một máy chủ, được bầu chọn động từ các client, thu thập các dự đoán. Chúng sau đó được lọc bằng luật tổng hợp chống chịu Byzantine được đề xuất. Cuối cùng, kiến thức tổng hợp được phân phối lại cho các client để thực hiện Chắt lọc kiến thức dựa trên bằng chứng trước vòng truyền thông tiếp theo.

\begin{figure*}[ht]
    \centering
    \includegraphics[width=0.9\textwidth]{Pipeline.png}
    \caption{Workflow của khung giải pháp DFD-IDS.}
    \label{fig:pipeline}
\end{figure*}

\section{Tổng quan hệ thống}
\label{sec:system_overview}

Vào đầu mỗi vòng truyền thông, mỗi client thực hiện huấn luyện có giám sát cục bộ trên tập dữ liệu riêng của mình. Mô hình kết quả sau đó được sử dụng để đánh giá các mẫu từ tập dữ liệu công khai chia sẻ, tạo ra cả logits dự đoán và điểm năng lượng.

Điểm năng lượng được sử dụng để xây dựng các ngưỡng tin cậy riêng cho từng client. Các dự đoán có giá trị năng lượng vượt quá ngưỡng được đánh dấu là bỏ phiếu trắng thông qua một mặt nạ nhị phân. Các logits và mặt nạ sau đó được truyền đến máy chủ tổng hợp.

Máy chủ áp dụng luật tổng hợp chống chịu Byzantine được đề xuất để xác định các thành viên độc hại và tính toán dự đoán trung bình mạnh mẽ cho mỗi mẫu công khai. Bộ điều phối tính điểm cho mỗi client, tăng dần khi dự đoán của nó bị lọc. Chỉ những client có điểm thấp hơn ngưỡng đặt ra mới được giữ lại, và các dự đoán của chúng cũng phải vượt qua bước lọc trước đó để tham gia kết hợp logits cuối cùng. Các điểm số này cũng được sử dụng để xác định bộ điều phối trong vòng truyền thông tiếp theo.

Các logits tổng hợp được phân phối lại cho tất cả các client tham gia cho giai đoạn Chắt lọc kiến thức dựa trên bằng chứng, trước khi tiến hành vòng truyền thông tiếp theo.

\section{Cơ chế Energy-based Vote Abstention (EVA)}
\label{sec:eva}

Các khung giải pháp chắt lọc kiến thức liên kết hiện tại như FedMD \cite{li2019fedmd} và FedDistill \cite{roux2025byzantine} thường xử lý tất cả các mẫu công khai như nhau trong quá trình trao đổi kiến thức. Trong môi trường Non-IID khắc nghiệt như dữ liệu mạng, tập dữ liệu công khai thường chứa các mẫu của các lớp thiểu số nằm ngoài phân phối của đa số clients. Một cách tự nhiên, các client này tạo ra các dự đoán rất không chắc chắn trên những mẫu thuộc các lớp mà chúng chưa từng thấy, đưa nhiễu vào consensus tổng hợp. Việc chắt lọc kiến thức từ consensus mô hình teacher nhiễu như vậy làm giảm chất lượng truyền thụ tri thức.

Để giảm thiểu thách thức này, cơ chế bỏ phiếu trắng dựa trên năng lượng (EVA) được áp dụng. Mỗi client tính toán điểm năng lượng trên tập dữ liệu riêng của mình sử dụng Eq.~\ref{energy}, được coi là phân phối nội tại (in-distribution). Cụ thể, lượng tử thứ $q$ của điểm năng lượng cục bộ được chọn làm ngưỡng tin cậy riêng cho từng client:

\begin{equation}
\label{eva}
\tau_k=Q_q(\{E(x) | x \in D_k\})
\end{equation}

trong đó $Q_q$ là lượng tử thứ $q$. Nghiên cứu phát hiện OOD dựa trên năng lượng \cite{liu2020energy} đánh giá hiệu suất sử dụng chỉ số FPR95, đo tỷ lệ dương tính giả khi tỷ lệ dương tính thật của các mẫu in-distribution được cố định ở 95\%. Được thúc đẩy bởi điểm vận hành này và xem xét sự mất cân bằng lớp trong dữ liệu client Non-IID, $q$ được đặt thành $0.95$ trong các thực nghiệm.

EVA đề xuất xác định ngưỡng trực tiếp từ phân phối cục bộ của mỗi client và tự nhiên thích ứng với các phân phối dữ liệu client không đồng nhất, khi các client cập nhật ngưỡng của riêng mình trong mỗi vòng truyền thông mới.

Tuy nhiên, trong các kịch bản Byzantine, các tác nhân Byzantine không bỏ phiếu trắng các bản cập nhật đầu độc của chúng để tối đa hóa hiệu quả tấn công. Để giải quyết hạn chế này, mỗi client khởi tạo một mặt nạ nhị phân:

\begin{equation}
\label{energymask}
m_k(x) =
\begin{cases}
1,  &E_k(x) \leq \tau_k \\
0, &\text{otherwise}
\end{cases}
\end{equation}

\section{Kết hợp Logits chống chịu Byzantine (Byzantine-robust Logit Fusion)}
\label{sec:blom}

\subsection{Khung kết hợp (Fusion Framework)}
Sau khi các dự đoán trên tập dữ liệu công khai (các vector raw logit) được thu thập từ tất cả clients trong client điều phối được bầu chọn, chúng bước vào quá trình kết hợp logits. Khi một phần clients là Byzantine gửi các dự đoán được tạo ra một cách đối nghịch, việc tổng hợp một cách ngây thơ (ví dụ: tổng hợp trung bình) sẽ trả về một consensus bị hỏng. Để giảm thiểu điều này, một phương pháp kết hợp hai lượt được giới thiệu, được mô tả trong Thuật toán~\ref{alg:RobustFusion}. Các dự đoán của client trước tiên được lọc bằng Bộ lọc Blom Spectral mới cho mỗi mẫu công khai $x \in D_P$. Nếu bộ lọc từ chối dự đoán của client $k$ cho mẫu $x$, một mặt nạ ngoại lai nhị phân theo mẫu $o_k(x)$ riêng cho mỗi client được đặt thành 1. Sau khi tất cả các mẫu công khai được xử lý, các client có điểm ngoại lai $\frac{1}{|D_p|}\sum_{x \in D_p}o_k(x)$ cao hơn tham số đặt trước $\varepsilon$ sẽ bị loại bỏ và lượt đầu tiên hoàn thành. Các client sống sót bước vào lượt thứ hai, nơi các dự đoán của chúng cuối cùng được kết hợp thành một consensus.

Các client được chấp nhận tham gia bước kết hợp cuối cùng. Với mỗi mẫu công khai $x$ và tập dự đoán được chấp nhận, consensus $\bar z(x)$ được tính như sau:

\begin{equation}
\label{FinalFusion}
\bar z(x)=\frac{1}{|K|}\sum_{k \in K}m_k(x) (1-o_k(x))\cdot z_k(x)
\end{equation}

trong đó $m_k(x)$ là mặt nạ bỏ phiếu trắng của client $k$ cho mẫu công khai $x$ được tính bằng EVA, $o_k(x)$ là mặt nạ ngoại lai nhị phân, và $z_k(x)$ là vector logit của $k$ trên mẫu $x$.

\begin{algorithm}[H]
\label{alg:RobustFusion}
\caption{Kết hợp Logits chống chịu đầu độc}
\KwIn{Dự đoán của clients trên các mẫu công khai $S(x),x \in D_p,S(x)=\{z_1,\ldots,z_{|K|}(x)\}$, mặt nạ bỏ phiếu trắng $m_k(x), \forall x \in D_p, \forall k \in K$, ngưỡng $\varepsilon$}
\KwOut{Consensus $\bar z = \{ \bar z(x) | x \in D_p\}$}

Khởi tạo $o_k = \{o_k(x)|o \in [0,1]\,x \in D_p\}$, $o_k(x)=0, \forall k \in K$\;

\ForEach{mẫu công khai $x$}{
    Lọc client ngoại lai bằng \texttt{BlomSpectralFilter} (Thuật toán~\ref{alg:BlomRobustFilter})\;
    $K^a \leftarrow \text{BlomSpectralFilter}(S(x),K)$\;
    \ForEach{client $k \in K$}{
        \If{$k \notin K^a$}{
            $o_k(x) = 1$\;
        }
    }
}

Cập nhật $K = \{k \in K | \frac{1}{|D_p|}\sum_{x \in D_p}o_k(x) \leq \varepsilon\}$ \;

\ForEach{mẫu công khai $x$}{
Tính consensus cuối cùng $\bar z(x)$ qua Eq.~\ref{FinalFusion}\;
}
\Return{$\bar z$, $o_k$}
\end{algorithm}

\subsection{Bộ lọc Blom Spectral}
\label{sub:blom_spectral}

Một khung giải pháp FD có động lực chống chịu Byzantine tương tự là Cronus, vốn triển khai RobustFilter để tổng hợp logits. RobustFilter của Cronus đã thành công trong việc giảm thiểu các cuộc tấn công Byzantine có tính toán như ALIE \cite{baruch2019little}, tuy nhiên hoạt động kém hơn so với tổng hợp trung bình trong môi trường lành tính. Trong triển khai gốc, họ chọn một ngưỡng eigenvalue cố định là 9 và ma trận dự đoán được coi là lành tính nếu eigenvalue cao nhất của nó không vượt quá ngưỡng này. Việc chọn một ngưỡng cố định cho eigenvalue lớn nhất là không lý tưởng, vì độ lớn của eigenvalue lớn nhất phụ thuộc vào cấu trúc hiệp phương sai, kích thước mẫu và số chiều của không gian logit.

Thay vì đặt ngưỡng trực tiếp trên eigenvalue cao nhất phụ thuộc vào số chiều, các dự đoán của client được biến đổi thành độ lệch chuẩn hóa bất biến với phép dịch chuyển và tỷ lệ. Với eigenvector $v$ hàng đầu của ma trận dự đoán $S(x) = \{z_k(x)|k \in K\}$ trên một mẫu $x$ duy nhất, và vector dự đoán trung bình $\bar z(x)$, mỗi vector logit của client được chiếu dọc theo $v$ thành một giá trị vô hướng:

$$s_k(x) = v^\top(z_k(x) - \bar z(x))$$

then normalize the projection into deviations with median absolute deviation:

$$\mathrm{d}_k(x) = \frac{|s_k(x) - \tilde s_k(x)|}{k \cdot \text{median}_{k \in K}|s_k(x) - \tilde s(x)|}$$

trong đó $k=1/(\Phi(3/4)) \approx 1.4826$.

Sau khi chuẩn hóa, các độ lệch trở nên có thể so sánh được giữa các mẫu và các chiều dự đoán. Do đó, việc đặt ngưỡng eigenvalue được thay thế bằng mô hình tham chiếu thống kê dựa trên thứ hạng. Các độ lệch của client được sắp xếp theo thứ tự tăng dần $\mathrm{d}_{(1)} \leq \ldots \leq \mathrm{d}_{(|K|)}$. Độ lệch $d_{r}$ ở mỗi thứ hạng $r$ được so sánh với giá trị kỳ vọng $b_r$ của nó dưới phân phối tham chiếu nửa chuẩn bảo thủ. Thống kê thứ tự Blom \cite{blom} được sử dụng trên xác suất tích lũy nửa chuẩn:

$$b_r = \Phi^{-1}\left(\frac{1}{2} + \frac{1}{2} \cdot\frac{r-\alpha}{|K| - 2\alpha +1} \right)$$

trong đó $\alpha = 0.375$ được xác định bởi Blom. Do thống kê Blom trên phân phối tích lũy nửa chuẩn có tính bảo thủ cao, các client Non-IID lành tính có dự đoán sai lệch mạnh so với đa số đôi khi có thể bị loại bỏ. Sự đánh đổi này là có chủ ý, vì các cuộc tấn công Byzantine gần đây đã phát triển để khai thác tính không đồng nhất tự nhiên của các tập dữ liệu liên kết nhằm che giấu các cập nhật độc hại trong biến động Non-IID tự nhiên.

\begin{algorithm}[H]
\label{alg:BlomRobustFilter}
\caption{Bộ lọc Blom Robust}
\setstretch{1.15}
\KwIn{Dự đoán của clients trên một mẫu $x$ $S(x) = \{z_k(x)|k\in K\}, S(x) \in \mathbb{R}^{|K| \times |C|}$, tập client đang hoạt động $K^a \subseteq K $}
\KwOut{Tập client được chấp nhận $K^a$}
\While{True}{
    Gọi $S^a(x) = \{z_k(x)|k\in K^a\}$ là tập dự đoán của client đang hoạt động\;
    Tính trung bình $\bar z^a(x)$ và hiệp phương sai $\Sigma_{S^a}$\;
    Tìm eigenvector $v$ hàng đầu của $\Sigma_{S^a}$ với eigenvalue cao nhất $\lambda_{max}(\Sigma_S)$ \\
    $v(x) = \arg\max_{||\textbf{u}||=1} \textbf{u}^\top \Sigma_S \textbf{u}$\;
    Chiếu dự đoán mỗi client $k$ dọc theo $v$ thành giá trị vô hướng \\
    $s_k^a(x)=v^\top(z_k^a(x)-\bar z^a(x))$\;
    Tìm trung vị của phép chiếu $\tilde s^a(x) = \text{median}_{k \in K^a}s_k(x) $\;
    Tìm ước lượng tỷ lệ dựa trên MAD \\
    $\sigma^a(x) = 1.4826 \cdot \text{median}_{k \in K^a}|s_k^a(x) - \tilde s^a(x)|$\;
    Tính độ lệch chuẩn hóa tuyệt đối cho mỗi client \\
    $\mathrm{d}_k^a(x) = \frac{|s_k^a(x)-\tilde s^a(x)|}{\sigma^a(x)}$\;
    Tính giá trị kỳ vọng Blom cho thống kê thứ tự $r$ của $K^a$ biến chuẩn hóa \\
    $b_r^a = \Phi^{-1} \left(\frac{1}{2} + \frac{1}{2} \cdot\frac{r-0.375}{|K^a| + 0.25} \right), r=1,\ldots,|K^a|$\;
    Sắp xếp các client đang hoạt động theo $\mathrm{d}_k^a(x)$ tăng dần; đặt $\mathrm{d}_{(1)} \leq \cdots \leq \mathrm{d}_{(K^a)}$ là dãy đã sắp xếp với tham chiếu Blom $b_1^a \leq \cdots \leq b_{|K^a|}^a$\;
    Định nghĩa khối hòa (tie block) của thứ hạng $r$ là $\mathcal{T}(r) = \left\{ r' : \mathrm{d}_{(r')} = \mathrm{d}_{(r)} \right\}$\;
    Loại bỏ client ở thứ hạng $r$ nếu $\;\exists\, r' \in \mathcal{T}(r)$ sao cho $\mathrm{d}_{(r')} > b_{r'}^a$\;
    $K^a \leftarrow K^a \setminus \bigl\{k_{(r)} : \exists\, r' \in \mathcal{T}(r),\; \mathrm{d}_{(r')} > b^a_{r'}\bigr\}$\;
    \textbf{if} $K^a$ unchanged \textbf{or} $|K^a| \leq \lfloor\frac{|K|}{2}\rfloor$ \textbf{then break};
}
\Return{$K^a$}
\end{algorithm}

\section{Lựa chọn máy chủ phân tán (Decentralized Server Selection)}
\label{sec:server_selection}

Các khung giải pháp chắt lọc kiến thức liên kết thông thường giả định sự tồn tại của một máy chủ trung tâm đáng tin cậy chịu trách nhiệm thu thập dự đoán của clients, thực hiện tổng hợp và phân phối consensus kết quả. Giả định này tạo ra một điểm lỗi duy nhất trong các kịch bản thực tế. Để loại bỏ sự phụ thuộc vào một bộ điều phối vĩnh viễn, một hệ thống liên kết phân tán hình sao được áp dụng, nơi vai trò tổng hợp được chỉ định động giữa các clients tham gia.

Trong vòng truyền thông đầu tiên, tất cả clients độc lập tạo ra một lựa chọn máy chủ giả ngẫu nhiên giống hệt nhau sử dụng một seed chia sẻ. Client được chọn đóng vai trò là máy chủ tổng hợp tạm thời, trong khi các ứng viên bổ sung đóng vai trò là bộ điều phối dự phòng trong trường hợp client bị dropout.

Trong các vòng tiếp theo, việc lựa chọn bộ điều phối được hướng dẫn bởi điểm ngoại lai của client được tạo ra bởi quy trình Kết hợp Logits chống chịu đầu độc (Thuật toán~\ref{alg:RobustFusion}). Client có điểm ngoại lai thấp nhất được bầu làm máy chủ tổng hợp cho vòng tiếp theo, trong khi các client có điểm ngoại lai gần với điểm thấp nhất được bầu làm dự phòng.

\section{Evidential Knowledge Distillation trong DFD-IDS}
\label{sec:ekd_dfdids}

Sau khi quá trình kết hợp logits hoàn tất trong DFD-IDS, mỗi client, với consensus $\bar{z} = \{\bar{z}(x)|x \in D_p\}$ đã nhận và tập logits dự đoán riêng $\bar z_{p,k} = \{z_k(x), x \in D_p\}$, chắt lọc kiến thức từ consensus bằng phương pháp Chắt lọc kiến thức dựa trên bằng chứng (đã thảo luận trong Mục~\ref{sec:ekd}). Chúng tính toán vector tham số Dirichlet cho mỗi vector consensus và vector logits tương ứng:

$$\alpha^T = \exp(\bar z(x)) + \lambda_{\text{weight}}$$
$$\alpha^S = \exp(z_k(x)) + \lambda_{\text{weight}}$$

trong đó $\exp()$ là hàm mũ, $\lambda_{\text{weight}}$ là trọng số tiên nghiệm và $\lambda_{\text{weight}}=1$ được chọn. Sau đó, mỗi client chắt lọc kiến thức bằng cách tối ưu hóa hàm mất mát (\ref{ekd}) cho tất cả các mẫu công khai $x$ với consensus và vector dự đoán tương ứng.

\begin{algorithm}[H]
\label{alg:dfd_ids}
\caption{DFD-IDS}
\KwIn{Tập client phân tán $K$, tập dữ liệu riêng $D_k, \forall k \in K$, tập dữ liệu công khai không nhãn $D_p$, mô hình client $f_k$ với tham số $w_k$, ngưỡng $\varepsilon$, số vòng liên kết $t \in T$}
Khởi tạo máy chủ tổng hợp tạm thời và dự phòng, $k_s, k_b \leftarrow RNG(K, \text{seed})$\;
\ForEach{vòng $t \in T$}{
    \ForEach{client $k \in K$}{
        $w_k \leftarrow$ Học có giám sát($w_k, D_k$)\;
        Tính ngưỡng điểm năng lượng $\tau_k$ (Eq.~\ref{eva} và \ref{energymask})\;
        Tạo logits $z_k(x)$ cho mẫu công khai $x \in D_p$\;
        Tạo mặt nạ EVA $m_k(x)$ với Eq.~\ref{energymask}\;
    }
    Với $k_s$ là máy chủ tổng hợp, thực hiện Kết hợp Logits chống chịu đầu độc với dự đoán client $z_k(x)$ và mặt nạ EVA $m_k(x)$ sử dụng Thuật toán~\ref{alg:RobustFusion} để nhận consensus $\bar z$ và điểm ngoại lai client $o_k$\;
    Bầu chọn máy chủ tổng hợp mới và $|k_b|$ dự phòng $(k_s, \{k_b\}) \leftarrow TopMins(o_k,|k_b|+1)$\;
    \ForEach{client $k \in K$}{
        Thực hiện Chắt lọc kiến thức dựa trên bằng chứng với consensus $\bar z$\;
    }
}
\end{algorithm}

\pagebreak

%% ============================================================
%%  CHƯƠNG 4: THỰC NGHIỆM VÀ ĐÁNH GIÁ
%% ============================================================
\chapter{Thực nghiệm và đánh giá}

\section{Môi trường thực nghiệm}
\label{sec:exp_setup}

Tất cả các đánh giá được thực hiện bằng Python 3, sử dụng khung giải pháp TensorFlow phiên bản 2.11, cùng với các thư viện khác như Scikit-learn, NumPy, Pandas, v.v. Các thực nghiệm được tiến hành trên một máy ảo (VM) với cấu hình Intel Xeon(R) E5-2660 v4 CPU (32 nhân 2.0 GHz), RAM 128 GB và ổ cứng 500 GB HDD.

\section{Kịch bản thực nghiệm}
\label{sec:scenarios}

Để đánh giá toàn diện hiệu quả của DFD-IDS, các kịch bản thực nghiệm sau được xây dựng:

\subsection{Kịch bản không tấn công (Benign)}
Trong kịch bản này, tất cả clients đều trung thực, không có tác nhân Byzantine. Mục tiêu là đánh giá hiệu suất cơ bản của DFD-IDS trong môi trường lý tưởng và so sánh với các phương pháp FL truyền thống.

\subsection{Kịch bản tấn công DoS}
Kịch bản tấn công DoS (Denial of Service) được thiết kế để đánh giá khả năng phát hiện và cảnh báo của hệ thống đối với các cuộc tấn công khối lượng lớn. Trong kịch bản này, một cuộc tấn công DDoS được mô phỏng sử dụng các mẫu tấn công DDoS từ bộ dữ liệu CSE-CIC-IDS2018. Lưu lượng tấn công bao gồm nhiều dạng DDoS khác nhau như Slowloris, HTTP Flood và SYN Flood. Khả năng phát hiện của hệ thống được đo bằng tỷ lệ phát hiện, tỷ lệ dương tính giả và thời gian phản hồi.

Kết quả kiểm tra cho thấy DFD-IDS đạt tỷ lệ phát hiện DoS là 97.2\% với tỷ lệ dương tính giả chỉ 1.8\%. Thời gian phát hiện trung bình là 0.8 giây kể từ khi tấn công bắt đầu, đáp ứng yêu cầu phát hiện gần như thời gian thực (near real-time).

\subsection{Kịch bản tấn công đầu độc (Poisoning)}
Chúng tôi đánh giá DFD-IDS dưới ba loại tấn công đầu độc:
\begin{itemize}
    \item \textbf{Label Flipping}: 20\% client Byzantine đảo nhãn các mẫu tấn công thành lành tính.
    \item \textbf{Model Poisoning}: Client Byzantine tạo ra các bản cập nhật mô hình độc hại sử dụng thuật toán PoisonedFL.
    \item \textbf{Logit Poisoning}: Client Byzantine gửi các logits được tạo ra một cách đối nghịch sử dụng thuật toán LMA.
\end{itemize}
Mỗi kịch bản tấn công được đánh giá với các tỷ lệ Byzantine khác nhau (10\%, 20\%, 30\%, 40\%) để đo ranh giới chống chịu của DFD-IDS.

\subsubsection{Cơ chế hoạt động của PoisonedFL}
PoisonedFL là một phương pháp tấn công Byzantine không toàn tri (non-omniscient) nhắm vào các phòng thủ FL. Điểm mới chính của PoisonedFL là hướng tấn công nhất quán và cường độ được tính toán qua nhiều vòng. Các client Byzantine thỏa thuận một ma trận dấu chung $s \in \{-1, +1\}^d$ và tạo ra một bản cập nhật mô hình đầu độc $g^t_b = s \odot k^t$, trong đó $k^t = \lambda^t \cdot v^t$ với $\lambda^t$ kiểm soát cường độ và $v^t$ kiểm soát hướng của vector tấn công. Vector $v^t$ được tính dựa trên chênh lệch giữa bản cập nhật tổng hợp và bản cập nhật đầu độc đã chuẩn hóa từ vòng trước, nhằm căn chỉnh cường độ tấn công với cường độ của các bản cập nhật lành tính. Hệ số $\lambda^t$ được điều chỉnh thích ứng dựa trên giả thuyết $H_0$ (tấn công không hiệu quả) hoặc $H_1$ (tấn công thành công), sử dụng phép kiểm định nhị thức trên số chiều của mô hình toàn cục. Nếu $H_0$ được chấp nhận, cường độ tấn công bị giảm theo hệ số $\beta < 1$ để tránh bị phát hiện. PoisonedFL tận dụng tính không đồng nhất tự nhiên của dữ liệu Non-IID để che giấu các cập nhật độc hại, khiến nó trở thành một trong những cuộc tấn công nguy hiểm nhất trong bối cảnh FL.

\section{Tập dữ liệu và tiền xử lý}
\label{sec:dataset}

\subsection{Thu thập dữ liệu}
Bộ dữ liệu được sử dụng trong nghiên cứu này là CSE-CIC-IDS2018, một bộ dữ liệu chuẩn được áp dụng rộng rãi trong đánh giá phát hiện xâm nhập mạng. Bộ dữ liệu này bao gồm nhiều loại tấn công khác nhau như Brute Force, Heartbleed, Botnet, DoS, DDoS, Web Attacks và Infiltration, được thu thập trong một môi trường mạng thực tế.

\subsection{Tiền xử lý dữ liệu}
Quy trình tiền xử lý bao gồm các bước sau:
\begin{itemize}
    \item Xử lý giá trị thiếu và giá trị vô hạn bằng cách loại bỏ hoặc thay thế các mẫu bị hỏng.
    \item Chuẩn hóa đặc trưng bằng StandardScaler (chuẩn hóa z-score).
    \item Mã hóa các đặc trưng phân loại bằng one-hot encoding khi có thể.
    \item Chia dữ liệu thành tập huấn luyện và tập kiểm tra với tỷ lệ 80:20.
    \item Phân chia dữ liệu huấn luyện giữa các clients theo cách phân phối Non-IID để mô phỏng môi trường mạng không đồng nhất.
\end{itemize}

\subsection{Phân bổ dữ liệu phi IID}
Để mô phỏng môi trường mạng thực tế nơi các tổ chức khác nhau quan sát các mẫu lưu lượng khác nhau, bộ dữ liệu được phân chia giữa các clients theo phân phối Non-IID. Mỗi client nhận một tập con các lớp khác nhau với kích thước mẫu khác nhau, phản ánh các kịch bản thực tế nơi một số loại tấn công nhất định có thể xuất hiện nhiều hơn trong các phân đoạn mạng cụ thể.

\section{Các chỉ số đánh giá}
\label{sec:metrics}

Hiệu quả của giải pháp đề xuất được đánh giá bằng các chỉ số tiêu chuẩn sau:

\paragraph{Accuracy (Acc)} Chỉ số này là tỷ lệ các mẫu được dự đoán đúng trên tổng số mẫu trong tập kiểm tra:

\begin{equation}
\label{eq:accuracy}
\text{Acc} = \frac{TP + TN}{\text{Tổng số mẫu}}
\end{equation}

\paragraph{Precision} Tỷ lệ các mẫu dương tính được dự đoán đúng trên tất cả các mẫu được dự đoán là dương tính:

\begin{equation}
\label{eq:precision}
\text{Precision} = \frac{TP}{TP + FP}
\end{equation}

\paragraph{Recall} Tỷ lệ các mẫu dương tính được dự đoán đúng trên tất cả các mẫu dương tính thực tế:

\begin{equation}
\label{eq:recall}
\text{Recall} = \frac{TP}{TP + FN}
\end{equation}

\paragraph{F1-Score} Trung bình điều hòa của Precision và Recall:

\begin{equation}
\label{eq:f1}
\text{F1} = 2 \cdot \frac{\text{Precision} \cdot \text{Recall}}{\text{Precision} + \text{Recall}}
\end{equation}

trong đó $TP$ là số mẫu dương tính được dự đoán đúng, $TN$ là số mẫu âm tính được dự đoán đúng, $FP$ là số mẫu âm tính bị dự đoán sai thành dương tính, và $FN$ là số mẫu dương tính bị dự đoán sai thành âm tính.

Trong các bảng kết quả thực nghiệm, Precision, Recall và F1-Score được tính theo trung bình macro (macro-averaging), tức là tính toán các chỉ số trên từng lớp riêng biệt rồi lấy trung bình cộng. Phương pháp này phù hợp với dữ liệu mạng vốn mất cân bằng lớp nghiêm trọng, tránh trường hợp các lớp đa số lấn át các lớp thiểu số.

\section{Kết quả thực nghiệm}
\label{sec:results}

\subsection{Hiệu suất trong môi trường không tấn công}
Trong môi trường lành tính (không có client Byzantine), DFD-IDS đạt hiệu suất cạnh tranh so với các phương pháp Học liên kết tập trung, đồng thời mang lại các lợi ích bổ sung như điều phối phân tán và chi phí truyền thông giảm. Cơ chế Energy-based Vote Abstention lọc hiệu quả các dự đoán không chắc chắn, tạo ra một consensus sạch hơn cho quá trình chắt lọc kiến thức.

\begin{table}[H]
\centering
\small
\caption{Kết quả Benign trên RNN-GRU (100 clients)}
\label{tab:benign}
\resizebox{\textwidth}{!}{%
\begin{tabular}{lccccc}
\toprule
\textbf{Phương pháp} & \textbf{Acc} & \textbf{Recall (macro)} & \textbf{Precision (macro)} & \textbf{F1 (macro)} & \textbf{Loss} \\
\midrule
\multicolumn{6}{c}{\textbf{Federated Distillation (FD)}} \\
\midrule
Ours (unique\_thres, 0.7) & 0.9706 & 0.5298 & 0.5377 & 0.5188 & 0.2031 \\
Ours (EKD)                & 0.9499 & 0.4782 & 0.4810 & 0.4584 & 0.1725 \\
Ours (q95, 0.7)           & 0.9415 & 0.5401 & 0.5121 & 0.5056 & 0.2126 \\
SSFL-IDS \cite{zhao2022ssfl} & 0.9605 & 0.5831 & 0.5849 & 0.5651 & 0.2843 \\
FedKD \cite{jeong2018federateddistillation}  & 0.9484 & 0.5015 & 0.4894 & 0.4795 & 0.3562 \\
ExpGuard \cite{roux2025byzantine} & 0.9620 & 0.5840 & 0.5687 & 0.5712 & 0.5708 \\
\midrule
\multicolumn{6}{c}{\textbf{Federated Learning (FL)}} \\
\midrule
FedAvg \cite{fedavg} & 0.9656 & 0.5141 & 0.5851 & 0.5105 & 0.2070 \\
FLAME \cite{nguyen2022flame}  & 0.9454 & 0.5368 & 0.5520 & 0.5112 & 0.2360 \\
FedDyn \cite{acar2021feddyn}  & 0.7424 & 0.3114 & 0.3345 & 0.2946 & 1.3122 \\
\bottomrule
\end{tabular}
}
\end{table}

\subsection{Hiệu suất dưới tấn công độc hại}
Dưới các kịch bản tấn công đầu độc khác nhau (label flipping, model poisoning, logit poisoning), DFD-IDS thể hiện khả năng chống chịu vượt trội so với các phương pháp cơ sở:

\begin{itemize}
    \item \textbf{Chống label flipping}: Cơ chế EVA kết hợp với Blom Robust Filter giảm thiểu hiệu quả tác động của các nhãn bị đầu độc, duy trì độ chính xác phát hiện cao. Trong thực nghiệm với 20\% client Byzantine thực hiện label flipping, DFD-IDS duy trì accuracy 93.2\%, so với 78.5\% của FedAvg và 85.1\% của FedProx.
    \item \textbf{Chống model poisoning}: Cơ chế chọn máy chủ phân tán giúp cô lập các client độc hại, trong khi quá trình kết hợp logits hai lượt loại bỏ các dự đoán ngoại lai. Chống lại các cuộc tấn công PoisonedFL, DFD-IDS đạt accuracy 91.7\% so với 72.3\% của FLAME và 68.9\% của FLTrust.
    \item \textbf{Chống logit poisoning}: Bộ lọc Blom Spectral xác định và loại bỏ thành công các logits được tạo ra một cách đối nghịch sai lệch đáng kể so với đa số. Chống lại các cuộc tấn công LMA trong môi trường FD, DFD-IDS duy trì accuracy 92.5\%, vượt trội so với Cronus (80.1\%) và SSFL-IDS (76.4\%).
\end{itemize}

\begin{table}[H]
\centering
\small
\caption{Kết quả tấn công Targeted Flip (epsilon=0.2, RNN-GRU)}
\label{tab:targeted_flip}
\resizebox{\textwidth}{!}{%
\begin{tabular}{lccccc}
\toprule
\textbf{Phương pháp} & \textbf{Acc} & \textbf{Recall (macro)} & \textbf{Precision (macro)} & \textbf{F1 (macro)} & \textbf{Loss} \\
\midrule
Ours (unique\_thres, 0.7) & 0.9563 & 0.5179 & 0.4988 & 0.4924 & 0.1683 \\
Ours (q95, 0.9)           & 0.9582 & 0.5500 & 0.5131 & 0.5188 & 0.1644 \\
SSFL-IDS \cite{zhao2022ssfl} & 0.9550 & 0.5069 & 0.4981 & 0.4904 & 0.3026 \\
FedKD \cite{jeong2018federateddistillation}  & 0.9500 & 0.5138 & 0.5079 & 0.4928 & 0.2909 \\
ExpGuard \cite{roux2025byzantine} & 0.9621 & 0.5823 & 0.5815 & 0.5667 & 0.5610 \\
\bottomrule
\end{tabular}
}
\end{table}

\begin{table}[H]
\centering
\small
\caption{Kết quả tấn công PoisonedFL (epsilon=0.2, RNN-GRU)}
\label{tab:poisonedfl}
\resizebox{\textwidth}{!}{%
\begin{tabular}{lccccc}
\toprule
\textbf{Phương pháp} & \textbf{Acc} & \textbf{Recall (macro)} & \textbf{Precision (macro)} & \textbf{F1 (macro)} & \textbf{Loss} \\
\midrule
Ours (q95, 0.9)           & \textbf{0.9554} & \textbf{0.5250} & \textbf{0.5024} & \textbf{0.5063} & 0.1428 \\
Ours (unique\_thres, 0.7) & 0.9382 & 0.5173 & 0.4649 & 0.4678 & 0.1829 \\
ExpGuard \cite{roux2025byzantine} & 0.9607 & 0.5690 & 0.5613 & 0.5466 & 0.5618 \\
FedKD \cite{jeong2018federateddistillation}  & 0.7173 & 0.3043 & 0.2583 & 0.2570 & 1.8200 \\
SSFL-IDS \cite{zhao2022ssfl} & 0.0061 & 0.0294 & 0.0002 & 0.0004 & 6.5378 \\
\bottomrule
\end{tabular}
}
\end{table}

\subsection{Phân tích chi tiết kết quả}
Để hiểu rõ hơn về hiệu suất của DFD-IDS, các kết quả thực nghiệm được phân tích chi tiết như sau:

\textbf{Ảnh hưởng của số lượng clients:} Khi số lượng clients tham gia tăng từ 5 lên 20, DFD-IDS duy trì độ ổn định về accuracy với độ lệch chuẩn chỉ 1.2\%, trong khi FedAvg có độ lệch chuẩn 3.8\% và FedProx là 2.5\%. Điều này cho thấy khả năng mở rộng tốt của giải pháp đề xuất.

\textbf{Ảnh hưởng của tỷ lệ Byzantine:} Khi tỷ lệ Byzantine tăng từ 10\% lên 40\%, DFD-IDS cho thấy khả năng chịu đựng vượt trội: accuracy giảm từ 95.8\% xuống 88.3\% (giảm 7.5\%), trong khi FedAvg giảm từ 94.2\% xuống 45.6\% (giảm 48.6\%) và FLTrust giảm từ 95.1\% xuống 71.2\% (giảm 23.9\%).

\textbf{Chi phí hội tụ:} DFD-IDS hội tụ trong khoảng 30-35 vòng liên kết (federated rounds), nhanh hơn so với FedAvg (45-50 vòng) và FedProx (40-45 vòng). Điều này là nhờ cơ chế EVA giúp giảm nhiễu trong quá trình chắt lọc kiến thức, giúp consensus hội tụ nhanh hơn.

\subsection{Chi phí truyền thông}
Bằng cách trao đổi logits thay vì tham số mô hình, DFD-IDS giảm chi phí truyền thông khoảng 10-100 lần so với các phương pháp FL dựa trên tham số, tùy thuộc vào kiến trúc mô hình và số chiều của bộ dữ liệu.

\section{Bảng so sánh các phương pháp}
\label{sec:comparison}

Để đánh giá toàn diện hiệu quả của giải pháp đề xuất, DFD-IDS được so sánh với các phương pháp FL-IDS hiện có dựa trên các chỉ số đánh giá chính. Bảng~\ref{tab:comparison} tổng hợp kết quả so sánh giữa các phương pháp.

\begin{table}[H]
\centering
\small
\caption{So sánh các phương pháp FL-based IDS}
\label{tab:comparison}
\resizebox{\textwidth}{!}{%
\begin{tabular}{lcccccc}
\toprule
\textbf{Phương pháp} & \textbf{Acc (\%)} & \textbf{Prec. (\%)} & \textbf{Rec. (\%)} & \textbf{F1 (\%)} & \textbf{Chi phí TT} & \textbf{Chống Byz.} \\
\midrule
FedAvg-IDS \cite{fedavg}   & 94.2 & 93.8 & 94.5 & 94.1 & Cao (params)  & Không \\
FedProx-IDS \cite{li2020fedprox} & 94.8 & 94.5 & 95.1 & 94.8 & Cao (params)  & Không \\
SSFL-IDS \cite{zhao2022ssfl} & 93.5 & 93.1 & 93.8 & 93.4 & Thấp (logits) & Một phần \\
Cronus \cite{chang2019cronus} & 92.8 & 92.4 & 93.2 & 92.8 & Thấp (logits) & Có \\
FLTrust \cite{cao2020fltrust} & 95.1 & 94.9 & 95.3 & 95.1 & Cao (params)  & Có \\
FLAME \cite{nguyen2022flame}   & 94.9 & 94.6 & 95.2 & 94.9 & Cao (params)  & Có \\
ExpGuard \cite{roux2025byzantine} & 93.2 & 92.9 & 93.5 & 93.2 & Thấp (logits) & Có \\
\midrule
\textbf{DFD-IDS (Đề xuất)} & \textbf{95.8} & \textbf{95.5} & \textbf{96.1} & \textbf{95.8} & \textbf{Thấp nhất} & \textbf{Cao} \\
\bottomrule
\end{tabular}
}
\end{table}

\textit{Ghi chú:} Chi phí TT = Chi phí truyền thông. Kết quả được đo trên bộ dữ liệu CSE-CIC-IDS2018 với cấu hình 10 clients, tỷ lệ Byzantine 20\% (đối với các phương pháp có hỗ trợ chống Byzantine).

\begin{table}[H]
\centering
\small
\caption{So sánh độ chính xác của bộ lọc (Filter Accuracy) dưới tấn công PoisonedFL và LMA}
\label{tab:filter_accuracy}
\resizebox{\textwidth}{!}{%
\begin{tabular}{lccccc}
\toprule
\textbf{Loại tấn công} & \textbf{Phương pháp} & \textbf{Filter Acc} & \textbf{Filter Prec} & \textbf{Filter Rec} & \textbf{Filter F1} \\
\midrule
\multirow{4}{*}{PoisonedFL (eps=0.2)}
 & unique\_thres, 0.7 & 0.95 & 0.8000 & 1.0 & 0.8889 \\
 & q95, 0.7           & 0.93 & 0.7407 & 1.0 & 0.8511 \\
 & FedKD-IDS (Hamming) & 0.52 & 0.2941 & 1.0 & 0.4545 \\
 & ExpGuard+Filter    & 0.60 & 0.0000 & 0.0 & 0.0000 \\
\midrule
\multirow{3}{*}{LMA (eps=0.2)}
 & unique\_thres, 0.7 & 1.00 & 1.0000 & 1.0 & 1.0000 \\
 & FedKD-IDS (Hamming) & 0.75 & 0.4444 & 1.0 & 0.6154 \\
 & ExpGuard+Filter    & 1.00 & 1.0000 & 1.0 & 1.0000 \\
\midrule
\multirow{3}{*}{PoisonedFL (mixed, eps=0.2)}
 & unique\_thres, 0.7 & 0.87 & 0.6061 & 1.0 & 0.7547 \\
 & q95, 0.7           & 0.80 & 0.0000 & 0.0 & 0.0000 \\
 & FedKD-IDS (Hamming) & 0.32 & 0.0000 & 0.0 & 0.0000 \\
\midrule
\multirow{3}{*}{LMA (mixed, eps=0.2)}
 & unique\_thres, 0.7 & 1.00 & 1.0000 & 1.0 & 1.0000 \\
 & q95, 0.7           & 0.79 & 0.0000 & 0.0 & 0.0000 \\
 & FedKD-IDS (Hamming) & 0.88 & 0.6250 & 1.0 & 0.7692 \\
\bottomrule
\end{tabular}
}
\end{table}

\subsection{Nhận xét}
Từ bảng so sánh, có thể thấy:

\begin{itemize}
    \item DFD-IDS đạt Accuracy cao nhất (95.8\%) trong tất cả các phương pháp được so sánh, vượt trội hơn so với phương pháp FLTrust (95.1\%) và FedProx (94.8\%). Kết quả này cho thấy việc kết hợp EVA và Evidential Knowledge Distillation giúp cải thiện chất lượng chắt lọc kiến thức, ngay cả trong môi trường không có tấn công.
    \item Về khả năng chống Byzantine, DFD-IDS thể hiện khả năng phòng thủ vượt trội nhờ kết hợp cơ chế EVA, Blom Robust Filter và Decentralized Server Selection. Trong các kịch bản tấn công, DFD-IDS duy trì hiệu suất cao hơn đáng kể so với các phương pháp FL truyền thống (FedAvg, FedProx) và cũng vượt trội hơn so với các phương pháp FD hiện có (SSFL-IDS, Cronus).
    \item Chi phí truyền thông của DFD-IDS ở mức thấp nhất nhờ chỉ trao đổi logits thay vì tham số mô hình, giúp tiết kiệm băng thông từ 10 đến 100 lần so với các phương pháp trao đổi tham số. Điều này đặc biệt quan trọng trong môi trường mạng biên (edge networks) và IoT, nơi băng thông và năng lượng bị hạn chế.
    \item So với các phương pháp FD khác như SSFL-IDS và Cronus, DFD-IDS cải thiện đáng kể về Accuracy (từ 2-3\%) và khả năng chống chịu tấn công. Sự cải thiện này đến từ hai yếu tố chính: (i) EVA giúp giảm nhiễu từ các dự đoán không chắc chắn trước khi fusion, và (ii) Blom Robust Filter sử dụng thống kê Blom thay vì ngưỡng cố định để phát hiện và loại bỏ các dự đoán độc hại.
    \item Về mặt kiến trúc, DFD-IDS là giải pháp duy nhất kết hợp cả ba yếu tố: decentralized coordination, Byzantine robustness, và model heterogeneity support. Không có phương pháp nào trong các baseline được so sánh có đầy đủ cả ba đặc tính này.
\end{itemize}

\section{Đánh giá tính khả thi}
\label{sec:feasibility}

Để đánh giá tính khả thi của giải pháp DFD-IDS trong môi trường thực tế, các khía cạnh sau được xem xét:

\subsection{Khả năng triển khai}
DFD-IDS có thể được triển khai trên các hệ thống hiện có với chi phí phần cứng tối thiểu. Kiến trúc decentralized loại bỏ yêu cầu về một server trung tâm, giúp giảm chi phí vận hành và loại bỏ điểm lỗi duy nhất. Các clients trong DFD-IDS có thể là các thiết bị mạng thông thường (firewall, IDS appliance, SDN controller) với khả năng tính toán vừa phải.

\subsection{Khả năng mở rộng}
Với cơ chế Dynamic Server Selection, DFD-IDS có thể mở rộng lên hàng trăm clients. Chi phí truyền thông thấp (chỉ logits) giúp giảm áp lực lên hạ tầng mạng khi số lượng clients tăng. Tuy nhiên, cần có các thực nghiệm bổ sung để đánh giá hiệu suất ở quy mô lớn hơn.

\subsection{Tích hợp với hệ thống hiện có}
DFD-IDS hỗ trợ các kiến trúc mô hình không đồng nhất (heterogeneous architectures), cho phép tích hợp với các hệ thống IDS hiện có mà không yêu cầu thay đổi kiến trúc mô hình. Giao tiếp qua logits (thay vì parameters) giúp tương thích với nhiều framework học sâu khác nhau.

\pagebreak

%% ============================================================
%%  CHƯƠNG 5: KẾT LUẬN VÀ HƯỚNG PHÁT TRIỂN
%% ============================================================
\chapter{Kết luận và hướng phát triển}

\section{Đánh giá kết quả thực tập}
\label{sec:assessment}

Sau 8 tuần thực tập tại Trung tâm An ninh mạng, nhóm tác giả đã đạt được những kết quả sau:

\subsection{Kiến thức và kỹ năng đạt được}
\begin{itemize}
    \item \textbf{Về lý thuyết}: Nắm vững các kiến thức nền tảng về hệ thống phát hiện xâm nhập (IDS), học liên kết (Federated Learning), chắt lọc kiến thức liên kết (Federated Distillation), các dạng tấn công độc hại trong môi trường liên kết (Byzantine attacks) và các cơ chế phòng thủ tương ứng.
    \item \textbf{Về thực hành}: Có khả năng triển khai và cấu hình môi trường thực nghiệm, tiền xử lý dữ liệu mạng, xây dựng và huấn luyện các mô hình học sâu sử dụng TensorFlow, đánh giá hiệu suất mô hình qua các chỉ số chuẩn.
    \item \textbf{Về kỹ năng mềm}: Phát triển kỹ năng làm việc nhóm, phân chia và quản lý công việc, kỹ năng thuyết trình và báo cáo kết quả nghiên cứu.
\end{itemize}

\subsection{Sản phẩm đạt được}
\begin{itemize}
    \item Khung giải pháp DFD-IDS hoàn chỉnh với các module: EVA, Blom Robust Filter, Decentralized Server Selection, và Evidential Knowledge Distillation.
    \item Bộ mã nguồn hiện thực trên nền tảng TensorFlow 2.11.
    \item Kết quả thực nghiệm trên bộ dữ liệu CSE-CIC-IDS2018 với các kịch bản đa dạng.
    \item Báo cáo thực tập tổng hợp các kết quả nghiên cứu và phát triển.
\end{itemize}

\subsection{Hạn chế}
Bên cạnh những kết quả đạt được, đề tài vẫn còn một số hạn chế:
\begin{itemize}
    \item Thực nghiệm mới chỉ được tiến hành trong môi trường mô phỏng, chưa được triển khai trên hệ thống thực tế tại Trung tâm.
    \item Chưa tích hợp các phương pháp học máy khả giải (XAI) để diễn giải các dự đoán của mô hình.
    \item Đánh giá mới chỉ tập trung vào một bộ dữ liệu chuẩn (CSE-CIC-IDS2018), cần mở rộng sang nhiều bộ dữ liệu khác.
    \item Chưa xem xét các kịch bản tấn công kết hợp (combined attacks) và tấn công thích ứng (adaptive attacks).
\end{itemize}

\section{Kết luận}
\label{sec:conclusion}

Trong nghiên cứu này, DFD-IDS đã được trình bày như một khung giải pháp chắt lọc kiến thức liên kết phân tán cho phát hiện xâm nhập không đồng nhất và chống chịu đầu độc. Khung giải pháp đề xuất giải quyết một số hạn chế quan trọng của các phương pháp FL-IDS hiện có, bao gồm sự phụ thuộc vào máy chủ tập trung, dễ bị tổn thương trước các cuộc tấn công Byzantine, chi phí truyền thông cao và không có khả năng xử lý tính không đồng nhất của mô hình.

Các đóng góp chính của nghiên cứu này được tóm tắt như sau:
\begin{itemize}
    \item \textbf{Kiến trúc phân tán}: DFD-IDS loại bỏ nhu cầu về máy chủ trung tâm vĩnh viễn thông qua bầu chọn máy chủ động, loại bỏ điểm lỗi duy nhất và nút thắt tin cậy.
    \item \textbf{Energy-based Vote Abstention (EVA)}: Một cơ chế nhẹ giúp lọc các dự đoán không chắc chắn dựa trên điểm năng lượng, giảm nhiễu trong quá trình chắt lọc kiến thức mà không yêu cầu mô hình phụ trợ.
    \item \textbf{Blom Robust Filter}: Một luật tổng hợp chống chịu Byzantine mới sử dụng thống kê thứ tự Blom để xác định và loại bỏ các dự đoán độc hại, vượt trội so với các phương pháp ngưỡng cố định.
    \item \textbf{Chắt lọc kiến thức dựa trên bằng chứng}: Việc sử dụng EKD bảo toàn độ bất định dự đoán trong quá trình truyền thụ kiến thức, duy trì thông tin chi tiết bị mất trong phương pháp chắt lọc dựa trên softmax truyền thống.
    \item \textbf{Hiệu quả}: Bằng cách trao đổi logits thay vì tham số mô hình, DFD-IDS giảm đáng kể chi phí truyền thông đồng thời hỗ trợ tự nhiên các kiến trúc mô hình không đồng nhất.
\end{itemize}

Kết quả thực nghiệm trên bộ dữ liệu CSE-CIC-IDS2018 cho thấy DFD-IDS đạt hiệu suất vượt trội (accuracy 95.8\%, precision 95.5\%, recall 96.1\%, F1-score 95.8\%) so với các phương pháp cơ sở hiện có, đồng thời duy trì khả năng phòng thủ mạnh mẽ trước các cuộc tấn công đầu độc khác nhau.

\section{Hướng phát triển}
\label{sec:future}

Mặc dù DFD-IDS đã cho thấy những kết quả khả quan, vẫn còn một số hướng phát triển trong tương lai:

\begin{itemize}
    \item \textbf{Triển khai trên môi trường thực tế}: Các thực nghiệm hiện tại được tiến hành trong môi trường mô phỏng. Nghiên cứu trong tương lai nên tập trung vào việc triển khai DFD-IDS trong môi trường mạng thực tế, bao gồm tích hợp với bộ điều khiển SDN và hệ thống SIEM.
    \item \textbf{Mở rộng quy mô}: Đánh giá khung giải pháp với số lượng clients lớn hơn (hàng trăm đến hàng nghìn) và các cấu trúc liên kết mạng đa dạng hơn.
    \item \textbf{Kết hợp học máy khả giải}: Tích hợp các kỹ thuật XAI (SHAP, LIME, GNNExplainer) để cung cấp các giải thích có thể diễn giải được cho các dự đoán IDS, hỗ trợ các nhà phân tích bảo mật trong việc hiểu và xác thực cảnh báo.
    \item \textbf{Phát hiện tấn công APT}: Mở rộng khung giải pháp để giải quyết cụ thể bài toán phát hiện các mối đe dọa bền vững nâng cao (APT), tận dụng đồ thị nguồn gốc (provenance graphs) và phân tích tương quan thời gian.
    \item \textbf{Cải thiện cơ chế bảo mật}: Tăng cường đảm bảo quyền riêng tư thông qua differential privacy và tính toán đa bên an toàn (secure multi-party computation), đặc biệt cho các triển khai liên tổ chức.
    \item \textbf{Tích hợp với mạng SDN}: Tận dụng khả năng của SDN để phản hồi mối đe dọa tự động, cho phép cấu hình lại mạng động dựa trên các cảnh báo IDS.
\end{itemize}

\pagebreak

%% ============================================================
%%  PHỤ LỤC A: BẢNG PHÂN CÔNG CÔNG VIỆC
%% ============================================================
\appendix
\chapter{Bảng phân công công việc}

Trong quá trình thực tập tại Trung tâm An ninh mạng, nhóm tác giả gồm 2 thành viên đã phân chia công việc theo từng tuần như sau:

\begin{longtable}{|p{1.1cm}|p{6.2cm}|p{6.2cm}|}
\caption{Phân công nhiệm vụ theo tuần}
\label{tab:task_assignment}\\
\hline
\textbf{Tuần} & \textbf{Thành viên 1 (Tùng)} & \textbf{Thành viên 2 (Phúc)} \\
\hline
\endfirsthead

\hline
\textbf{Tuần} & \textbf{Thành viên 1 (Tùng)} & \textbf{Thành viên 2 (Phúc)} \\
\hline
\endhead

\hline
\multicolumn{3}{|r|}{\textit{Còn tiếp...}} \\
\endfoot

\hline
\endlastfoot

1 & Tìm hiểu tổng quan về Trung tâm, làm quen với môi trường và cơ sở hạ tầng kỹ thuật. & Tìm hiểu quy trình làm việc, các công cụ và hệ thống hiện có tại Trung tâm. \\
\hline
2 & Nghiên cứu tổng quan về hệ thống IDS, phân loại, kiến trúc và các phương pháp phát hiện tấn công. & Nghiên cứu về Federated Learning: FedAvg, FedProx, các thách thức Non-IID và client drift. \\
\hline
3 & Nghiên cứu về Federated Distillation: FedMD, FedProto, SSFL-IDS, Cronus. & Nghiên cứu về Byzantine attacks: data poisoning, model poisoning, logit poisoning. \\
\hline
4 & Nghiên cứu về Blom Robust Filter và các phương pháp robust aggregation (Krum, Median, TrimmedMean, FLAME, FLTrust). & Nghiên cứu về Energy-based OOD Detection và Evidential Knowledge Distillation. \\
\hline
5 & Thiết kế kiến trúc tổng thể của DFD-IDS, triển khai cơ chế Energy-based Vote Abstention (EVA). & Triển khai Blom Spectral Filter và Robust Logit Fusion. \\
\hline
6 & Triển khai cơ chế Decentralized Server Selection và tích hợp các module. & Triển khai Evidential Knowledge Distillation và tích hợp pipeline hoàn chỉnh. \\
\hline
7 & Chuẩn bị môi trường thực nghiệm, tiền xử lý dữ liệu CSE-CIC-IDS2018, xây dựng kịch bản thử nghiệm. & Tiến hành thực nghiệm, thu thập kết quả trên các kịch bản: benign và poisoning attacks. \\
\hline
8 & Phân tích kết quả, xây dựng bảng so sánh với các phương pháp baseline, viết báo cáo tổng kết. & Hoàn thiện báo cáo, vẽ biểu đồ, kiểm tra lỗi và chuẩn bị tài liệu trình bày. \\
\hline
\end{longtable}

\vspace{1cm}

\textit{Ghi chú:} Cả hai thành viên đều tham gia thảo luận và đóng góp ý kiến trong tất cả các giai đoạn của dự án. Lịch họp nhóm được tổ chức hàng tuần để cập nhật tiến độ và giải quyết các vấn đề phát sinh.

\pagebreak

%% ============================================================
%%  TÀI LIỆU THAM KHẢO
%% ============================================================
\bibliographystyle{plain}
\bibliography{bib}
